% ================================================================
% O singură sursă, două PDF-uri (compilează fiecare de două ori):
%
%   La clasă (interactiv, soluția se dezvăluie pe bucăți):
%     pdflatex -jobname=lesson-interactive lesson.tex
%
%   Compact (toți pașii într-un singur slide, pentru elevi / README):
%     pdflatex -jobname=lesson "\PassOptionsToClass{handout}{beamer}% ================================================================
% O singură sursă, două PDF-uri (compilează fiecare de două ori):
%
%   La clasă (interactiv, soluția se dezvăluie pe bucăți):
%     pdflatex -jobname=lesson-interactive lesson.tex
%
%   Compact (toți pașii într-un singur slide, pentru elevi / README):
%     pdflatex -jobname=lesson "\PassOptionsToClass{handout}{beamer}% ================================================================
% O singură sursă, două PDF-uri (compilează fiecare de două ori):
%
%   La clasă (interactiv, soluția se dezvăluie pe bucăți):
%     pdflatex -jobname=lesson-interactive lesson.tex
%
%   Compact (toți pașii într-un singur slide, pentru elevi / README):
%     pdflatex -jobname=lesson "\PassOptionsToClass{handout}{beamer}% ================================================================
% O singură sursă, două PDF-uri (compilează fiecare de două ori):
%
%   La clasă (interactiv, soluția se dezvăluie pe bucăți):
%     pdflatex -jobname=lesson-interactive lesson.tex
%
%   Compact (toți pașii într-un singur slide, pentru elevi / README):
%     pdflatex -jobname=lesson "\PassOptionsToClass{handout}{beamer}\input{lesson.tex}"
% ================================================================
\documentclass[9pt]{beamer}
\usepackage{theme-sorting}

% Mod interactiv: întrebarea apare prima, apoi soluția se dezvăluie pe bucăți
% (\item<2->, \uncover<3->{...}, \begin{remarca}<4->{Concluzie}); hint-ul se
% folosește doar unde elevii chiar trebuie împinși să se gândească.
% Totul se afișează cumulativ: ce apare la pasul 2 rămâne vizibil după.
\newenvironment<>{intrebare}{%
  \begin{uncoverenv}#1%
  \begin{block}{Întrebare}}{\end{block}\end{uncoverenv}}
% Hint și remarcă: text inline cu etichetă emerald, fără cutie
\newenvironment<>{hint}{%
  \begin{uncoverenv}#1%
  \par\textcolor{igMarine}{\textbf{Hint:}}\ }{\par\end{uncoverenv}}
\newenvironment<>{remarca}[1]{%
  \begin{uncoverenv}#2%
  \par\smallskip\textcolor{igMarine}{\textbf{#1:}}\ }{\par\end{uncoverenv}}

% Celulă de tabel vizibilă de la pasul #1
\newcommand{\at}[2]{\uncover<#1->{#2}}

% Diagrame cu tablouri: celulă normală (N), comparată / interschimbată (H),
% la locul definitiv (S), minim sau element inserat (M)
\tikzset{
  cN/.style={draw=igNavy!55,fill=white,minimum size=0.62cm,inner sep=0pt,font=\small},
  cH/.style={cN,fill=igBlue!22,draw=igBlue,thick},
  cS/.style={cN,fill=igSlate!26},
  cM/.style={cN,fill=igOchre!32,draw=igOchre,thick}
}
% \arow{y}{valoare/stil, ...}{text}  (5 elemente)
\newcommand{\arow}[3]{%
  \foreach \v/\t [count=\i from 0] in {#2} {\node[c\t] at (\i*0.68,#1) {\v};}%
  \node[anchor=west,font=\footnotesize] at (5*0.68-0.25,#1) {#3};}

\title{Algoritmi de sortare}
\subtitle{Bubble $\cdot$ Selection $\cdot$ Insertion $\cdot$ STL $\cdot$ Counting Sort}
\author[A. M. Ciobanu]{Andrei Mihai Ciobanu}
\institute[Algolymp]{Algolymp Summer School}
\date{}
\titlegraphic{\includegraphics[width=3cm]{../../../resources/algolymp-logo.png}}

% ================================================================
\begin{document}

\igtitleslide

% ================================================================
\section{Problema sortării}

\begin{frame}{Rolul sortării}
  Sortarea este una dintre cele mai frecvente operații în programare.
  \begin{itemize}
    \item Căutarea rapidă a unui element necesită un șir sortat (tehnică studiată în lecția următoare)
    \item Multe probleme devin simple după sortare: al $k$-lea cel mai mare element, perechi cu suma dată, eliminarea duplicatelor
    \item Date reale: clasamente, liste de note, ierarhii
  \end{itemize}
  \begin{block}{Problema sortării}
    Dat un vector $v[0..n-1]$, să se rearanjeze elementele astfel încât
    $v[0] \leq v[1] \leq \ldots \leq v[n-1]$.
  \end{block}
  \begin{intrebare}<2->
    Cum se verifică dacă un șir de $n$ numere conține două elemente egale, dacă șirul este nesortat, respectiv sortat?
  \end{intrebare}
  \uncover<3->{Nesortat: se compară toate perechile, deci circa $\frac{n(n-1)}{2}$ comparații. Sortat: elementele egale sunt adiacente, deci este suficientă o singură parcurgere, cu $n-1$ comparații.}
\end{frame}

\begin{frame}{Inversiuni}
  O \textbf{inversiune} este o pereche de poziții $i < j$ cu $v[i] > v[j]$.
  \begin{intrebare}
    Câte inversiuni are șirul $[5,\ 3,\ 8,\ 1,\ 4]$?
  \end{intrebare}
  \begin{hint}<2->
    Se numără, pentru fiecare element, valorile mai mici aflate în dreapta lui.
  \end{hint}
  \begin{itemize}
    \item<3-> $5$: $(5,3),\ (5,1),\ (5,4)$ \quad $\rightarrow$ 3 inversiuni
    \item<3-> $3$: $(3,1)$ \quad $\rightarrow$ 1 inversiune
    \item<3-> $8$: $(8,1),\ (8,4)$ \quad $\rightarrow$ 2 inversiuni
    \item<3-> $1$ și $4$: nicio inversiune \quad $\Rightarrow$ \textbf{6 inversiuni} în total
  \end{itemize}
  \begin{remarca}<4->{Legătura cu sortarea}
    Un șir este sortat dacă și numai dacă nu are inversiuni. Interschimbarea a două elemente \textbf{adiacente} aflate în ordine greșită elimină exact o inversiune, iar numărul maxim de inversiuni este $\frac{n(n-1)}{2}$ (șir descrescător).
  \end{remarca}
\end{frame}

% ================================================================
\section{Bubble Sort}

\begin{frame}{Bubble Sort \textemdash{} idee}
  Se compară perechile de elemente adiacente, iar elementele aflate în ordine greșită se interschimbă.
  \begin{intrebare}
    Ce se întâmplă cu cel mai mare element după o parcurgere completă, de la stânga la dreapta, a șirului $[5,\ 3,\ 8,\ 1,\ 4]$?
  \end{intrebare}
  \begin{center}
  \begin{tikzpicture}
    \uncover<2->{\arow{0}{5/N,3/N,8/N,1/N,4/N}{șirul inițial}}
    \uncover<3->{\arow{-0.7}{3/H,5/H,8/N,1/N,4/N}{$5 > 3$: se interschimbă}}
    \uncover<4->{\arow{-1.4}{3/N,5/H,8/H,1/N,4/N}{$5 < 8$: rămân pe loc}}
    \uncover<5->{\arow{-2.1}{3/N,5/N,1/H,8/H,4/N}{$8 > 1$: se interschimbă}}
    \uncover<6->{\arow{-2.8}{3/N,5/N,1/N,4/H,8/H}{$8 > 4$: se interschimbă}}
    \uncover<7->{\arow{-3.5}{3/N,5/N,1/N,4/N,8/S}{sfârșitul pasului 1}}
  \end{tikzpicture}
  \end{center}
  \begin{remarca}<7->{Concluzie}
    După pasul $1$, cel mai mare element se află pe ultima poziție, la locul său definitiv.
  \end{remarca}
  \beginnote{Vizualizare pas cu pas: \href{https://visualgo.net/en/sorting}{visualgo.net/en/sorting}}
\end{frame}

\begin{frame}{Bubble Sort \textemdash{} pașii următori}
  După pasul $i$, ultimele $i$ elemente sunt la locul lor definitiv, deci pasul $i+1$ nu mai ajunge la ele.
  \begin{center}
  \begin{tikzpicture}
    \arow{0}{3/N,5/N,1/N,4/N,8/S}{după pasul 1 \ (3 interschimbări)}
    \uncover<2->{\arow{-0.7}{3/N,1/N,4/N,5/S,8/S}{după pasul 2 \ (2 interschimbări)}}
    \uncover<3->{\arow{-1.4}{1/S,3/S,4/S,5/S,8/S}{după pasul 3 \ (1 interschimbare)}}
  \end{tikzpicture}
  \end{center}
  \begin{intrebare}<4->
    Câte interschimbări a făcut algoritmul în total și ce legătură are numărul cu inversiunile șirului inițial?
  \end{intrebare}
  \uncover<5->{$3 + 2 + 1 = 6$ interschimbări, adică exact cele $6$ inversiuni ale șirului $[5, 3, 8, 1, 4]$.}
  \begin{itemize}
    \item<6-> Pasul $4$ nu mai modifică nimic: pentru $n$ elemente sunt necesari cel mult $n-1$ pași
  \end{itemize}
\end{frame}

\begin{frame}[fragile]{Bubble Sort \textemdash{} implementare}
  \begin{lstlisting}
void bubbleSort(int v[], int n) {
    for (int i = 0; i < n - 1; i++) {
        for (int j = 0; j < n - 1 - i; j++) {
            if (v[j] > v[j + 1]) {
                swap(v[j], v[j + 1]);
            }
        }
    }
}
  \end{lstlisting}
  \begin{itemize}
    \item Bucla exterioară: $n-1$ pași
    \item Bucla interioară: la pasul $i$ se fac $n-1-i$ comparații
    \item Total comparații: $(n-1) + (n-2) + \ldots + 1 = \frac{n(n-1)}{2}$
  \end{itemize}
  \begin{intrebare}<2->
    De ce bucla interioară se oprește la $n-1-i$, nu la $n-1$?
  \end{intrebare}
  \uncover<3->{Ultimele $i$ elemente sunt deja la locul lor, deci nu mai este nevoie să fie comparate.}
  \beginnote{swap \textemdash{} funcție din STL, interschimbă valorile a două variabile}
\end{frame}

\begin{frame}[fragile]{Bubble Sort \textemdash{} optimizare și complexitate}
  \begin{intrebare}
    Ce complexitate are algoritmul de mai sus pe un șir deja sortat?
  \end{intrebare}
  \uncover<2->{
    $O(n^2)$: comparațiile se fac oricum. Algoritmul se poate opri imediat ce o parcurgere completă nu mai găsește nicio inversare.}
  \begin{uncoverenv}<3->
    \begin{lstlisting}
void bubbleSort(int v[], int n) {
    bool ok = true;            // exista inversiuni
    while (ok) {
        ok = false;            // presupunem ca nu mai sunt
        for (int j = 0; j < n - 1; j++) {
            if (v[j] > v[j + 1]) {
                swap(v[j], v[j + 1]);
                ok = true;     // am gasit o inversiune
            }
        }
    }
}
    \end{lstlisting}
  \end{uncoverenv}
  \begin{itemize}
    \item<4-> \textbf{Best case} (șir sortat): $O(n)$ \textemdash{} o singură parcurgere, nicio inversare
    \item<4-> \textbf{Worst case} (șir sortat invers) și \textbf{average case}: $O(n^2)$
  \end{itemize}
\end{frame}

% ================================================================
\section{Selection Sort}

\begin{frame}{Selection Sort \textemdash{} idee}
  La fiecare pas se găsește minimul din partea nesortată și se aduce pe prima poziție a acesteia.
  \begin{intrebare}
    Cum se sortează șirul $[5,\ 3,\ 8,\ 1,\ 4]$ prin alegerea repetată a minimului?
  \end{intrebare}
  \begin{center}
  \begin{tikzpicture}
    \uncover<2->{\arow{0}{5/H,3/N,8/N,1/M,4/N}{minim $1$ $\rightarrow$ se interschimbă cu poziția 0}}
    \uncover<3->{\arow{-0.7}{1/S,3/M,8/N,5/N,4/N}{minim $3$, deja la locul lui}}
    \uncover<4->{\arow{-1.4}{1/S,3/S,8/H,5/N,4/M}{minim $4$ $\rightarrow$ se interschimbă cu poziția 2}}
    \uncover<5->{\arow{-2.1}{1/S,3/S,4/S,5/M,8/N}{minim $5$, deja la locul lui}}
    \uncover<6->{\arow{-2.8}{1/S,3/S,4/S,5/S,8/S}{șir sortat}}
  \end{tikzpicture}
  \end{center}
  \begin{remarca}<6->{Concluzie}
    După pasul $i$, primele $i+1$ elemente sunt la locul lor definitiv.
  \end{remarca}
  \beginnote{Vizualizare pas cu pas: \href{https://visualgo.net/en/sorting}{visualgo.net/en/sorting}}
\end{frame}

\begin{frame}[fragile]{Selection Sort \textemdash{} implementare}
  \begin{lstlisting}
void selectionSort(int v[], int n) {
    for (int i = 0; i < n - 1; i++) {
        int pozMin = i;                  // pozitia minimului din v[i..n-1]
        for (int j = i + 1; j < n; j++) {
            if (v[j] < v[pozMin]) {
                pozMin = j;
            }
        }
        swap(v[i], v[pozMin]);
    }
}
  \end{lstlisting}
  \begin{itemize}
    \item \textbf{Best, worst și average case}: $O(n^2)$ \textemdash{} mereu se fac $\frac{n(n-1)}{2}$ comparații
    \item Se fac cel mult $n-1$ interschimbări (câte una pe pas), mai puține decât la Bubble Sort
  \end{itemize}
  \begin{intrebare}<2->
    Ce complexitate are algoritmul pe un șir deja sortat?
  \end{intrebare}
  \uncover<3->{Tot $O(n^2)$: minimul se caută oricum în toată partea nesortată.}
  \beginnote{swap \textemdash{} funcție din STL, interschimbă valorile a două variabile}
\end{frame}

% ================================================================
\section{Insertion Sort}

\begin{frame}{Insertion Sort \textemdash{} idee}
  Se construiește un prefix sortat, iar fiecare element nou este inserat la locul potrivit în acest prefix.
  \begin{intrebare}
    Cum se extinde un prefix sortat cu încă un element, fără a sorta din nou tot prefixul?
  \end{intrebare}
  \begin{center}
  \begin{tikzpicture}
    \uncover<2->{\arow{0}{5/S,3/M,8/N,1/N,4/N}{se inserează $3$ $\rightarrow$ $[3, 5]$}}
    \uncover<3->{\arow{-0.7}{3/S,5/S,8/M,1/N,4/N}{se inserează $8$ \textemdash{} rămâne pe loc}}
    \uncover<4->{\arow{-1.4}{3/S,5/S,8/S,1/M,4/N}{se inserează $1$ $\rightarrow$ ajunge la început}}
    \uncover<5->{\arow{-2.1}{1/S,3/S,5/S,8/S,4/M}{se inserează $4$ $\rightarrow$ între $3$ și $5$}}
    \uncover<6->{\arow{-2.8}{1/S,3/S,4/S,5/S,8/S}{șir sortat}}
  \end{tikzpicture}
  \end{center}
  \begin{remarca}<6->{Concluzie}
    Elementul inserat este mutat spre stânga cât timp este mai mic decât vecinul său.
  \end{remarca}
  \beginnote{Vizualizare pas cu pas: \href{https://visualgo.net/en/sorting}{visualgo.net/en/sorting}}
\end{frame}

\begin{frame}[fragile]{Insertion Sort \textemdash{} implementare}
  \begin{lstlisting}
void insertionSort(int v[], int n) {
    for (int i = 1; i < n; i++) {
        int j = i;
        while (j > 0 && v[j - 1] > v[j]) {
            swap(v[j - 1], v[j]);
            j--;
        }
    }
}
  \end{lstlisting}
  \begin{itemize}
    \item \textbf{Best case} (șir sortat): $O(n)$ \textemdash{} nicio deplasare
    \item \textbf{Worst case} (șir sortat invers) și \textbf{average case}: $O(n^2)$
  \end{itemize}
  \begin{intrebare}<2->
    De ce condiția buclei conține $j > 0$?
  \end{intrebare}
  \uncover<3->{Fără ea, pentru $j = 0$ s-ar accesa $v[-1]$, adică în afara tabloului.}
  \beginnote{swap \textemdash{} funcție din STL, interschimbă valorile a două variabile}
\end{frame}

\begin{frame}{Cele trei metode pătratice}
  \begin{intrebare}
    Care dintre cele trei metode face cele mai puține interschimbări și care se comportă bine pe un șir aproape sortat?
  \end{intrebare}
  \uncover<2->{
  \begin{center}
  \begin{tabular}{lccc}
    \toprule
    Algoritm & Best case & Worst case & Interschimbări \\
    \midrule
    Bubble Sort (optimizat) & $O(n)$   & $O(n^2)$ & \at{2}{egal cu numărul de inversiuni} \\
    Selection Sort          & $O(n^2)$ & $O(n^2)$ & \at{2}{cel mult $n-1$} \\
    Insertion Sort          & $O(n)$   & $O(n^2)$ & \at{2}{egal cu numărul de inversiuni} \\
    \bottomrule
  \end{tabular}
  \end{center}}
  \begin{itemize}
    \item<3-> \textbf{Selection Sort} face cele mai puține interschimbări
    \item<3-> \textbf{Insertion Sort} este eficient pe șiruri aproape sortate: timpul depinde de numărul de inversiuni
  \end{itemize}
\end{frame}

% ================================================================
\section{STL sort}

\begin{frame}[fragile]{STL sort \textemdash{} utilizare de bază}
  C++ oferă funcția \lstinline|sort| din header-ul \lstinline|<algorithm>|, cu complexitate $O(n \log n)$.
  \begin{intrebare}
    Cum se sortează un tablou $v$ cu $n$ elemente crescător, respectiv descrescător?
  \end{intrebare}
  \begin{uncoverenv}<2->
    \begin{lstlisting}
#include <algorithm>
using namespace std;

int v[] = {5, 3, 8, 1, 4};
int n = 5;

sort(v, v + n);                    // crescator: 1 3 4 5 8
sort(v, v + n, greater<int>());    // descrescator: 8 5 4 3 1
sort(v + l, v + r + 1);            // doar pozitiile l..r

vector<int> w = {5, 3, 8, 1, 4};
sort(w.begin(), w.end());          // pentru vector
    \end{lstlisting}
  \end{uncoverenv}
  \beginnote{vector, begin, end, greater \textemdash{} din STL}
\end{frame}

\begin{frame}[fragile]{STL sort \textemdash{} comparatori}
  Criteriul de sortare poate fi schimbat printr-o \textbf{funcție de comparare}, care întoarce \lstinline|true| dacă primul argument trebuie să apară înaintea celui de-al doilea.
  \begin{intrebare}
    Cum se sortează numerele după valoarea absolută?
  \end{intrebare}
  \begin{uncoverenv}<2->
    \begin{lstlisting}
bool cmpAbs(int a, int b) {
    return abs(a) < abs(b);
}

sort(v, v + n, cmpAbs);
    \end{lstlisting}
  \end{uncoverenv}
  \begin{uncoverenv}<3->
    Pentru perechi: crescător după primul element, iar la egalitate descrescător după al doilea.
    \begin{lstlisting}
bool cmpPair(pair<int, int> a, pair<int, int> b) {
    if (a.first != b.first)
        return a.first < b.first;
    return a.second > b.second;
}

sort(pv, pv + m, cmpPair);   // pv: tablou de pair<int, int>
    \end{lstlisting}
  \end{uncoverenv}
  \beginnote{pair, abs \textemdash{} din STL}
\end{frame}

\begin{frame}[fragile]{STL sort \textemdash{} reguli pentru comparatori}
  \begin{intrebare}
    Ce poate face \lstinline|sort| dacă un comparator întoarce \lstinline|true| și pentru două elemente egale, de exemplu \lstinline|return a <= b;|?
  \end{intrebare}
  \uncover<2->{Comportament nedefinit: programul poate da rezultate greșite sau se poate opri cu eroare (acces în afara tabloului).}
  \begin{block}<3->{Reguli pentru un comparator corect}
    \begin{itemize}
      \item \lstinline|cmp(a, a)| trebuie să întoarcă \lstinline|false|: se folosesc \lstinline|<| și \lstinline|>|, niciodată \lstinline|<=| sau \lstinline|>=|
      \item dacă \lstinline|cmp(a, b)| este \lstinline|true|, atunci \lstinline|cmp(b, a)| trebuie să fie \lstinline|false|
      \item relația trebuie să fie tranzitivă: din \lstinline|cmp(a, b)| și \lstinline|cmp(b, c)| rezultă \lstinline|cmp(a, c)|
    \end{itemize}
  \end{block}
  \begin{itemize}
    \item<4-> \lstinline|sort| nu garantează ordinea relativă a elementelor egale; \lstinline|stable_sort| o păstrează (sortare \textbf{stabilă})
  \end{itemize}
\end{frame}

% ================================================================
\section{Counting Sort}

\begin{frame}{Counting Sort \textemdash{} idee}
  Elementele nu se compară între ele: se numără de câte ori apare fiecare valoare, apoi șirul se reconstruiește.
  \begin{intrebare}
    Cum se poate sorta șirul $v = [3,\ 1,\ 4,\ 1,\ 5,\ 3,\ 3]$ cu valori în intervalul $[0, 5]$, fără nicio comparație între elemente?
  \end{intrebare}
  \begin{hint}<2->
    Valorile sunt numere mici: pot fi folosite ca indici într-un tablou de contoare.
  \end{hint}
  \uncover<3->{
    \begin{center}
    \begin{tabular}{ccccccc}
      \toprule
      Valoare & 0 & 1 & 2 & 3 & 4 & 5 \\
      \midrule
      Contor & 0 & \at{4}{2} & 0 & \at{4}{3} & \at{4}{1} & \at{4}{1} \\
      \bottomrule
    \end{tabular}
    \end{center}}
  \begin{remarca}<5->{Rezultat}
    Fiecare valoare se scrie de câte ori apare: $[1, 1, 3, 3, 3, 4, 5]$.
  \end{remarca}
  \beginnote{Vizualizare pas cu pas: \href{https://visualgo.net/en/sorting}{visualgo.net/en/sorting}}
\end{frame}

\begin{frame}[fragile]{Counting Sort \textemdash{} implementare}
  \begin{lstlisting}
const int MAXV = 1000000;
int cnt[MAXV + 1];                    // contoarele, initial 0

void countingSort(int v[], int n) {
    int maxVal = 0;
    for (int i = 0; i < n; i++) {
        cnt[v[i]]++;                  // numaram aparitiile
        if (v[i] > maxVal) maxVal = v[i];
    }
    int idx = 0;                      // reconstruim sirul sortat
    for (int val = 0; val <= maxVal; val++) {
        for (int i = 0; i < cnt[val]; i++) {
            v[idx++] = val;
        }
        cnt[val] = 0;                 // contoarele raman 0 pentru apelul urmator
    }
}
  \end{lstlisting}
  \begin{itemize}
    \item Complexitate $O(n + k)$: $n$ pentru numărare, $k$ pentru parcurgerea valorilor posibile
  \end{itemize}
\end{frame}

\begin{frame}{Counting Sort \textemdash{} limitări}
  \begin{intrebare}
    În ce situații Counting Sort nu este potrivit?
  \end{intrebare}
  \begin{block}<2->{Când funcționează}
    \begin{itemize}
      \item Valori întregi nenegative
      \item $k$ mic (de exemplu note $1$\textendash{}$10$, vârste $0$\textendash{}$150$, valori $\leq 10^6$)
    \end{itemize}
  \end{block}
  \begin{block}<3->{Când nu funcționează}
    \begin{itemize}
      \item Valori negative (se adună o constantă tuturor valorilor) sau reale (trebuie transformate mai întâi)
      \item $k$ foarte mare: valori până la $10^9$ ar necesita un tablou de aproximativ $4$ GB
    \end{itemize}
  \end{block}
\end{frame}

% ================================================================
\section{Comparație}

\begin{frame}{Comparația algoritmilor de sortare}
  \begin{center}
  \begin{tabular}{llll}
    \toprule
    Algoritm & Timp & Spațiu suplimentar & Condiții \\
    \midrule
    Bubble Sort    & $O(n^2)$      & $O(1)$      & orice tip comparabil \\
    Selection Sort & $O(n^2)$      & $O(1)$      & orice tip comparabil \\
    Insertion Sort & $O(n^2)$      & $O(1)$      & orice tip comparabil \\
    STL sort       & $O(n \log n)$ & $O(\log n)$ & orice tip comparabil \\
    Counting Sort  & $O(n + k)$    & $O(k)$      & întregi în $[0, k]$ \\
    \bottomrule
  \end{tabular}
  \end{center}
  \begin{intrebare}<2->
    Ce algoritm este potrivit pentru $n = 10^5$ valori până la $10^9$?
  \end{intrebare}
  \uncover<3->{STL sort: metodele pătratice fac circa $10^{10}$ operații, iar Counting Sort ar necesita un tablou de $10^9$ elemente.}
  \begin{itemize}
    \item<4-> În practică se folosește \textbf{STL sort}, dacă nu există restricții speciale
    \item<4-> Counting Sort este mai rapid decât STL sort când $k = O(n)$
    \item<4-> Bubble Sort este util didactic, dar se evită în concursuri
  \end{itemize}
\end{frame}

% ================================================================
\section{Recapitulare}

\begin{frame}[fragile]{Recapitulare}
  \begin{center}
  \begin{tabular}{ll}
    \toprule
    Operație & Cod C++ \\
    \midrule
    Sortare crescătoare    & \lstinline|sort(v, v + n)| \\
    Sortare descrescătoare & \lstinline|sort(v, v + n, greater<int>())| \\
    Sortare a pozițiilor $l..r$ & \lstinline[keepspaces=true]|sort(v + l, v + r + 1)| \\
    Comparator propriu     & \lstinline|sort(v, v + n, cmp)| \\
    \bottomrule
  \end{tabular}
  \end{center}
  \begin{itemize}
    \item \textbf{Bubble, Selection, Insertion Sort}: $O(n^2)$; Bubble și Insertion fac exact atâtea interschimbări câte inversiuni există
    \item \textbf{STL sort}: $O(n \log n)$, soluția implicită; comparatorul nu trebuie să întoarcă \lstinline|true| pentru elemente egale
    \item \textbf{Counting Sort}: $O(n + k)$, doar pentru întregi mici
  \end{itemize}
\end{frame}

\begin{frame}{Probleme de antrenament}
  \begin{block}{Ușoară}
    \begin{enumerate}
      \item \href{https://code.algolymp.com/problems/1081}{\textbf{Even Sort}}
            \textcolor{gray}{(AlgoCode \#1081)}
    \end{enumerate}
  \end{block}
  \begin{block}{Medii}
    \begin{enumerate}
      \item \href{https://code.algolymp.com/problems/1077}{\textbf{K Sort}}
            \textcolor{gray}{(AlgoCode \#1077)}
      \item \href{https://code.algolymp.com/problems/1076}{\textbf{Sort Primes}}
            \textcolor{gray}{(AlgoCode \#1076)}
    \end{enumerate}
  \end{block}
  \begin{block}{Grea}
    \begin{enumerate}
      \item \href{https://code.algolymp.com/problems/1082}{\textbf{Array Task}}
            \textcolor{gray}{(AlgoCode \#1082)}
    \end{enumerate}
  \end{block}
\end{frame}

\begin{frame}{Resurse}
  \begin{block}{Materiale}
    \begin{enumerate}
      \item \href{https://www.pbinfo.ro/articole/5587/sortarea-tablourilor}{\textbf{pbinfo.ro}}
            \textemdash{} Sortarea tablourilor (RO)
      \item \href{https://www.pbinfo.ro/articole/5563/sortarea-prin-selectie}{\textbf{pbinfo.ro}}
            \textemdash{} Sortare prin selecție (RO)
      \item \href{https://www.pbinfo.ro/articole/23702/standard-template-library-stl}{\textbf{pbinfo.ro}}
            \textemdash{} STL în română
      \item \href{https://cp-algorithms.com/sequences/bubble-sort.html}{\textbf{cp-algorithms.com}}
            \textemdash{} Bubble Sort (EN)
      \item \href{https://en.cppreference.com/w/cpp/algorithm/sort}{\textbf{cppreference.com}}
            \textemdash{} std::sort (EN)
      \item \href{https://en.wikipedia.org/wiki/Quicksort}{\textbf{wikipedia.org}}
            \textemdash{} Quicksort și \href{https://en.wikipedia.org/wiki/Merge_sort}{Merge Sort}, algoritmi $O(n \log n)$ (EN)
    \end{enumerate}
  \end{block}
  \begin{block}{Vizualizări}
    \begin{enumerate}
      \item \href{https://visualgo.net/en/sorting}{\textbf{VisuAlgo}}
            \textemdash{} animații pas cu pas pentru Bubble, Selection, Insertion, Merge, Quick, Counting și Radix Sort, cu date introduse de utilizator (EN)
      \item \href{https://www.toptal.com/developers/sorting-algorithms}{\textbf{Toptal}}
            \textemdash{} opt algoritmi comparați pe date aleatoare, aproape sortate, inversate și cu puține valori distincte (EN)
      \item \href{https://www.cs.usfca.edu/~galles/visualization/ComparisonSort.html}{\textbf{USFCA}}
            \textemdash{} Comparison Sorting Visualization (EN)
      \item \href{https://panthema.net/2013/sound-of-sorting/}{\textbf{Sound of Sorting}}
            \textemdash{} vizualizare și sunet pentru operațiile de sortare (EN)
      \item \href{https://www.youtube.com/watch?v=kPRA0W1kECg}{\textbf{YouTube}}
            \textemdash{} 15 Sorting Algorithms in 6 Minutes (EN)
    \end{enumerate}
  \end{block}
\end{frame}

\end{document}
"
% ================================================================
\documentclass[9pt]{beamer}
\usepackage{theme-sorting}

% Mod interactiv: întrebarea apare prima, apoi soluția se dezvăluie pe bucăți
% (\item<2->, \uncover<3->{...}, \begin{remarca}<4->{Concluzie}); hint-ul se
% folosește doar unde elevii chiar trebuie împinși să se gândească.
% Totul se afișează cumulativ: ce apare la pasul 2 rămâne vizibil după.
\newenvironment<>{intrebare}{%
  \begin{uncoverenv}#1%
  \begin{block}{Întrebare}}{\end{block}\end{uncoverenv}}
% Hint și remarcă: text inline cu etichetă emerald, fără cutie
\newenvironment<>{hint}{%
  \begin{uncoverenv}#1%
  \par\textcolor{igMarine}{\textbf{Hint:}}\ }{\par\end{uncoverenv}}
\newenvironment<>{remarca}[1]{%
  \begin{uncoverenv}#2%
  \par\smallskip\textcolor{igMarine}{\textbf{#1:}}\ }{\par\end{uncoverenv}}

% Celulă de tabel vizibilă de la pasul #1
\newcommand{\at}[2]{\uncover<#1->{#2}}

% Diagrame cu tablouri: celulă normală (N), comparată / interschimbată (H),
% la locul definitiv (S), minim sau element inserat (M)
\tikzset{
  cN/.style={draw=igNavy!55,fill=white,minimum size=0.62cm,inner sep=0pt,font=\small},
  cH/.style={cN,fill=igBlue!22,draw=igBlue,thick},
  cS/.style={cN,fill=igSlate!26},
  cM/.style={cN,fill=igOchre!32,draw=igOchre,thick}
}
% \arow{y}{valoare/stil, ...}{text}  (5 elemente)
\newcommand{\arow}[3]{%
  \foreach \v/\t [count=\i from 0] in {#2} {\node[c\t] at (\i*0.68,#1) {\v};}%
  \node[anchor=west,font=\footnotesize] at (5*0.68-0.25,#1) {#3};}

\title{Algoritmi de sortare}
\subtitle{Bubble $\cdot$ Selection $\cdot$ Insertion $\cdot$ STL $\cdot$ Counting Sort}
\author[A. M. Ciobanu]{Andrei Mihai Ciobanu}
\institute[Algolymp]{Algolymp Summer School}
\date{}
\titlegraphic{\includegraphics[width=3cm]{../../../resources/algolymp-logo.png}}

% ================================================================
\begin{document}

\igtitleslide

% ================================================================
\section{Problema sortării}

\begin{frame}{Rolul sortării}
  Sortarea este una dintre cele mai frecvente operații în programare.
  \begin{itemize}
    \item Căutarea rapidă a unui element necesită un șir sortat (tehnică studiată în lecția următoare)
    \item Multe probleme devin simple după sortare: al $k$-lea cel mai mare element, perechi cu suma dată, eliminarea duplicatelor
    \item Date reale: clasamente, liste de note, ierarhii
  \end{itemize}
  \begin{block}{Problema sortării}
    Dat un vector $v[0..n-1]$, să se rearanjeze elementele astfel încât
    $v[0] \leq v[1] \leq \ldots \leq v[n-1]$.
  \end{block}
  \begin{intrebare}<2->
    Cum se verifică dacă un șir de $n$ numere conține două elemente egale, dacă șirul este nesortat, respectiv sortat?
  \end{intrebare}
  \uncover<3->{Nesortat: se compară toate perechile, deci circa $\frac{n(n-1)}{2}$ comparații. Sortat: elementele egale sunt adiacente, deci este suficientă o singură parcurgere, cu $n-1$ comparații.}
\end{frame}

\begin{frame}{Inversiuni}
  O \textbf{inversiune} este o pereche de poziții $i < j$ cu $v[i] > v[j]$.
  \begin{intrebare}
    Câte inversiuni are șirul $[5,\ 3,\ 8,\ 1,\ 4]$?
  \end{intrebare}
  \begin{hint}<2->
    Se numără, pentru fiecare element, valorile mai mici aflate în dreapta lui.
  \end{hint}
  \begin{itemize}
    \item<3-> $5$: $(5,3),\ (5,1),\ (5,4)$ \quad $\rightarrow$ 3 inversiuni
    \item<3-> $3$: $(3,1)$ \quad $\rightarrow$ 1 inversiune
    \item<3-> $8$: $(8,1),\ (8,4)$ \quad $\rightarrow$ 2 inversiuni
    \item<3-> $1$ și $4$: nicio inversiune \quad $\Rightarrow$ \textbf{6 inversiuni} în total
  \end{itemize}
  \begin{remarca}<4->{Legătura cu sortarea}
    Un șir este sortat dacă și numai dacă nu are inversiuni. Interschimbarea a două elemente \textbf{adiacente} aflate în ordine greșită elimină exact o inversiune, iar numărul maxim de inversiuni este $\frac{n(n-1)}{2}$ (șir descrescător).
  \end{remarca}
\end{frame}

% ================================================================
\section{Bubble Sort}

\begin{frame}{Bubble Sort \textemdash{} idee}
  Se compară perechile de elemente adiacente, iar elementele aflate în ordine greșită se interschimbă.
  \begin{intrebare}
    Ce se întâmplă cu cel mai mare element după o parcurgere completă, de la stânga la dreapta, a șirului $[5,\ 3,\ 8,\ 1,\ 4]$?
  \end{intrebare}
  \begin{center}
  \begin{tikzpicture}
    \uncover<2->{\arow{0}{5/N,3/N,8/N,1/N,4/N}{șirul inițial}}
    \uncover<3->{\arow{-0.7}{3/H,5/H,8/N,1/N,4/N}{$5 > 3$: se interschimbă}}
    \uncover<4->{\arow{-1.4}{3/N,5/H,8/H,1/N,4/N}{$5 < 8$: rămân pe loc}}
    \uncover<5->{\arow{-2.1}{3/N,5/N,1/H,8/H,4/N}{$8 > 1$: se interschimbă}}
    \uncover<6->{\arow{-2.8}{3/N,5/N,1/N,4/H,8/H}{$8 > 4$: se interschimbă}}
    \uncover<7->{\arow{-3.5}{3/N,5/N,1/N,4/N,8/S}{sfârșitul pasului 1}}
  \end{tikzpicture}
  \end{center}
  \begin{remarca}<7->{Concluzie}
    După pasul $1$, cel mai mare element se află pe ultima poziție, la locul său definitiv.
  \end{remarca}
  \beginnote{Vizualizare pas cu pas: \href{https://visualgo.net/en/sorting}{visualgo.net/en/sorting}}
\end{frame}

\begin{frame}{Bubble Sort \textemdash{} pașii următori}
  După pasul $i$, ultimele $i$ elemente sunt la locul lor definitiv, deci pasul $i+1$ nu mai ajunge la ele.
  \begin{center}
  \begin{tikzpicture}
    \arow{0}{3/N,5/N,1/N,4/N,8/S}{după pasul 1 \ (3 interschimbări)}
    \uncover<2->{\arow{-0.7}{3/N,1/N,4/N,5/S,8/S}{după pasul 2 \ (2 interschimbări)}}
    \uncover<3->{\arow{-1.4}{1/S,3/S,4/S,5/S,8/S}{după pasul 3 \ (1 interschimbare)}}
  \end{tikzpicture}
  \end{center}
  \begin{intrebare}<4->
    Câte interschimbări a făcut algoritmul în total și ce legătură are numărul cu inversiunile șirului inițial?
  \end{intrebare}
  \uncover<5->{$3 + 2 + 1 = 6$ interschimbări, adică exact cele $6$ inversiuni ale șirului $[5, 3, 8, 1, 4]$.}
  \begin{itemize}
    \item<6-> Pasul $4$ nu mai modifică nimic: pentru $n$ elemente sunt necesari cel mult $n-1$ pași
  \end{itemize}
\end{frame}

\begin{frame}[fragile]{Bubble Sort \textemdash{} implementare}
  \begin{lstlisting}
void bubbleSort(int v[], int n) {
    for (int i = 0; i < n - 1; i++) {
        for (int j = 0; j < n - 1 - i; j++) {
            if (v[j] > v[j + 1]) {
                swap(v[j], v[j + 1]);
            }
        }
    }
}
  \end{lstlisting}
  \begin{itemize}
    \item Bucla exterioară: $n-1$ pași
    \item Bucla interioară: la pasul $i$ se fac $n-1-i$ comparații
    \item Total comparații: $(n-1) + (n-2) + \ldots + 1 = \frac{n(n-1)}{2}$
  \end{itemize}
  \begin{intrebare}<2->
    De ce bucla interioară se oprește la $n-1-i$, nu la $n-1$?
  \end{intrebare}
  \uncover<3->{Ultimele $i$ elemente sunt deja la locul lor, deci nu mai este nevoie să fie comparate.}
  \beginnote{swap \textemdash{} funcție din STL, interschimbă valorile a două variabile}
\end{frame}

\begin{frame}[fragile]{Bubble Sort \textemdash{} optimizare și complexitate}
  \begin{intrebare}
    Ce complexitate are algoritmul de mai sus pe un șir deja sortat?
  \end{intrebare}
  \uncover<2->{
    $O(n^2)$: comparațiile se fac oricum. Algoritmul se poate opri imediat ce o parcurgere completă nu mai găsește nicio inversare.}
  \begin{uncoverenv}<3->
    \begin{lstlisting}
void bubbleSort(int v[], int n) {
    bool ok = true;            // exista inversiuni
    while (ok) {
        ok = false;            // presupunem ca nu mai sunt
        for (int j = 0; j < n - 1; j++) {
            if (v[j] > v[j + 1]) {
                swap(v[j], v[j + 1]);
                ok = true;     // am gasit o inversiune
            }
        }
    }
}
    \end{lstlisting}
  \end{uncoverenv}
  \begin{itemize}
    \item<4-> \textbf{Best case} (șir sortat): $O(n)$ \textemdash{} o singură parcurgere, nicio inversare
    \item<4-> \textbf{Worst case} (șir sortat invers) și \textbf{average case}: $O(n^2)$
  \end{itemize}
\end{frame}

% ================================================================
\section{Selection Sort}

\begin{frame}{Selection Sort \textemdash{} idee}
  La fiecare pas se găsește minimul din partea nesortată și se aduce pe prima poziție a acesteia.
  \begin{intrebare}
    Cum se sortează șirul $[5,\ 3,\ 8,\ 1,\ 4]$ prin alegerea repetată a minimului?
  \end{intrebare}
  \begin{center}
  \begin{tikzpicture}
    \uncover<2->{\arow{0}{5/H,3/N,8/N,1/M,4/N}{minim $1$ $\rightarrow$ se interschimbă cu poziția 0}}
    \uncover<3->{\arow{-0.7}{1/S,3/M,8/N,5/N,4/N}{minim $3$, deja la locul lui}}
    \uncover<4->{\arow{-1.4}{1/S,3/S,8/H,5/N,4/M}{minim $4$ $\rightarrow$ se interschimbă cu poziția 2}}
    \uncover<5->{\arow{-2.1}{1/S,3/S,4/S,5/M,8/N}{minim $5$, deja la locul lui}}
    \uncover<6->{\arow{-2.8}{1/S,3/S,4/S,5/S,8/S}{șir sortat}}
  \end{tikzpicture}
  \end{center}
  \begin{remarca}<6->{Concluzie}
    După pasul $i$, primele $i+1$ elemente sunt la locul lor definitiv.
  \end{remarca}
  \beginnote{Vizualizare pas cu pas: \href{https://visualgo.net/en/sorting}{visualgo.net/en/sorting}}
\end{frame}

\begin{frame}[fragile]{Selection Sort \textemdash{} implementare}
  \begin{lstlisting}
void selectionSort(int v[], int n) {
    for (int i = 0; i < n - 1; i++) {
        int pozMin = i;                  // pozitia minimului din v[i..n-1]
        for (int j = i + 1; j < n; j++) {
            if (v[j] < v[pozMin]) {
                pozMin = j;
            }
        }
        swap(v[i], v[pozMin]);
    }
}
  \end{lstlisting}
  \begin{itemize}
    \item \textbf{Best, worst și average case}: $O(n^2)$ \textemdash{} mereu se fac $\frac{n(n-1)}{2}$ comparații
    \item Se fac cel mult $n-1$ interschimbări (câte una pe pas), mai puține decât la Bubble Sort
  \end{itemize}
  \begin{intrebare}<2->
    Ce complexitate are algoritmul pe un șir deja sortat?
  \end{intrebare}
  \uncover<3->{Tot $O(n^2)$: minimul se caută oricum în toată partea nesortată.}
  \beginnote{swap \textemdash{} funcție din STL, interschimbă valorile a două variabile}
\end{frame}

% ================================================================
\section{Insertion Sort}

\begin{frame}{Insertion Sort \textemdash{} idee}
  Se construiește un prefix sortat, iar fiecare element nou este inserat la locul potrivit în acest prefix.
  \begin{intrebare}
    Cum se extinde un prefix sortat cu încă un element, fără a sorta din nou tot prefixul?
  \end{intrebare}
  \begin{center}
  \begin{tikzpicture}
    \uncover<2->{\arow{0}{5/S,3/M,8/N,1/N,4/N}{se inserează $3$ $\rightarrow$ $[3, 5]$}}
    \uncover<3->{\arow{-0.7}{3/S,5/S,8/M,1/N,4/N}{se inserează $8$ \textemdash{} rămâne pe loc}}
    \uncover<4->{\arow{-1.4}{3/S,5/S,8/S,1/M,4/N}{se inserează $1$ $\rightarrow$ ajunge la început}}
    \uncover<5->{\arow{-2.1}{1/S,3/S,5/S,8/S,4/M}{se inserează $4$ $\rightarrow$ între $3$ și $5$}}
    \uncover<6->{\arow{-2.8}{1/S,3/S,4/S,5/S,8/S}{șir sortat}}
  \end{tikzpicture}
  \end{center}
  \begin{remarca}<6->{Concluzie}
    Elementul inserat este mutat spre stânga cât timp este mai mic decât vecinul său.
  \end{remarca}
  \beginnote{Vizualizare pas cu pas: \href{https://visualgo.net/en/sorting}{visualgo.net/en/sorting}}
\end{frame}

\begin{frame}[fragile]{Insertion Sort \textemdash{} implementare}
  \begin{lstlisting}
void insertionSort(int v[], int n) {
    for (int i = 1; i < n; i++) {
        int j = i;
        while (j > 0 && v[j - 1] > v[j]) {
            swap(v[j - 1], v[j]);
            j--;
        }
    }
}
  \end{lstlisting}
  \begin{itemize}
    \item \textbf{Best case} (șir sortat): $O(n)$ \textemdash{} nicio deplasare
    \item \textbf{Worst case} (șir sortat invers) și \textbf{average case}: $O(n^2)$
  \end{itemize}
  \begin{intrebare}<2->
    De ce condiția buclei conține $j > 0$?
  \end{intrebare}
  \uncover<3->{Fără ea, pentru $j = 0$ s-ar accesa $v[-1]$, adică în afara tabloului.}
  \beginnote{swap \textemdash{} funcție din STL, interschimbă valorile a două variabile}
\end{frame}

\begin{frame}{Cele trei metode pătratice}
  \begin{intrebare}
    Care dintre cele trei metode face cele mai puține interschimbări și care se comportă bine pe un șir aproape sortat?
  \end{intrebare}
  \uncover<2->{
  \begin{center}
  \begin{tabular}{lccc}
    \toprule
    Algoritm & Best case & Worst case & Interschimbări \\
    \midrule
    Bubble Sort (optimizat) & $O(n)$   & $O(n^2)$ & \at{2}{egal cu numărul de inversiuni} \\
    Selection Sort          & $O(n^2)$ & $O(n^2)$ & \at{2}{cel mult $n-1$} \\
    Insertion Sort          & $O(n)$   & $O(n^2)$ & \at{2}{egal cu numărul de inversiuni} \\
    \bottomrule
  \end{tabular}
  \end{center}}
  \begin{itemize}
    \item<3-> \textbf{Selection Sort} face cele mai puține interschimbări
    \item<3-> \textbf{Insertion Sort} este eficient pe șiruri aproape sortate: timpul depinde de numărul de inversiuni
  \end{itemize}
\end{frame}

% ================================================================
\section{STL sort}

\begin{frame}[fragile]{STL sort \textemdash{} utilizare de bază}
  C++ oferă funcția \lstinline|sort| din header-ul \lstinline|<algorithm>|, cu complexitate $O(n \log n)$.
  \begin{intrebare}
    Cum se sortează un tablou $v$ cu $n$ elemente crescător, respectiv descrescător?
  \end{intrebare}
  \begin{uncoverenv}<2->
    \begin{lstlisting}
#include <algorithm>
using namespace std;

int v[] = {5, 3, 8, 1, 4};
int n = 5;

sort(v, v + n);                    // crescator: 1 3 4 5 8
sort(v, v + n, greater<int>());    // descrescator: 8 5 4 3 1
sort(v + l, v + r + 1);            // doar pozitiile l..r

vector<int> w = {5, 3, 8, 1, 4};
sort(w.begin(), w.end());          // pentru vector
    \end{lstlisting}
  \end{uncoverenv}
  \beginnote{vector, begin, end, greater \textemdash{} din STL}
\end{frame}

\begin{frame}[fragile]{STL sort \textemdash{} comparatori}
  Criteriul de sortare poate fi schimbat printr-o \textbf{funcție de comparare}, care întoarce \lstinline|true| dacă primul argument trebuie să apară înaintea celui de-al doilea.
  \begin{intrebare}
    Cum se sortează numerele după valoarea absolută?
  \end{intrebare}
  \begin{uncoverenv}<2->
    \begin{lstlisting}
bool cmpAbs(int a, int b) {
    return abs(a) < abs(b);
}

sort(v, v + n, cmpAbs);
    \end{lstlisting}
  \end{uncoverenv}
  \begin{uncoverenv}<3->
    Pentru perechi: crescător după primul element, iar la egalitate descrescător după al doilea.
    \begin{lstlisting}
bool cmpPair(pair<int, int> a, pair<int, int> b) {
    if (a.first != b.first)
        return a.first < b.first;
    return a.second > b.second;
}

sort(pv, pv + m, cmpPair);   // pv: tablou de pair<int, int>
    \end{lstlisting}
  \end{uncoverenv}
  \beginnote{pair, abs \textemdash{} din STL}
\end{frame}

\begin{frame}[fragile]{STL sort \textemdash{} reguli pentru comparatori}
  \begin{intrebare}
    Ce poate face \lstinline|sort| dacă un comparator întoarce \lstinline|true| și pentru două elemente egale, de exemplu \lstinline|return a <= b;|?
  \end{intrebare}
  \uncover<2->{Comportament nedefinit: programul poate da rezultate greșite sau se poate opri cu eroare (acces în afara tabloului).}
  \begin{block}<3->{Reguli pentru un comparator corect}
    \begin{itemize}
      \item \lstinline|cmp(a, a)| trebuie să întoarcă \lstinline|false|: se folosesc \lstinline|<| și \lstinline|>|, niciodată \lstinline|<=| sau \lstinline|>=|
      \item dacă \lstinline|cmp(a, b)| este \lstinline|true|, atunci \lstinline|cmp(b, a)| trebuie să fie \lstinline|false|
      \item relația trebuie să fie tranzitivă: din \lstinline|cmp(a, b)| și \lstinline|cmp(b, c)| rezultă \lstinline|cmp(a, c)|
    \end{itemize}
  \end{block}
  \begin{itemize}
    \item<4-> \lstinline|sort| nu garantează ordinea relativă a elementelor egale; \lstinline|stable_sort| o păstrează (sortare \textbf{stabilă})
  \end{itemize}
\end{frame}

% ================================================================
\section{Counting Sort}

\begin{frame}{Counting Sort \textemdash{} idee}
  Elementele nu se compară între ele: se numără de câte ori apare fiecare valoare, apoi șirul se reconstruiește.
  \begin{intrebare}
    Cum se poate sorta șirul $v = [3,\ 1,\ 4,\ 1,\ 5,\ 3,\ 3]$ cu valori în intervalul $[0, 5]$, fără nicio comparație între elemente?
  \end{intrebare}
  \begin{hint}<2->
    Valorile sunt numere mici: pot fi folosite ca indici într-un tablou de contoare.
  \end{hint}
  \uncover<3->{
    \begin{center}
    \begin{tabular}{ccccccc}
      \toprule
      Valoare & 0 & 1 & 2 & 3 & 4 & 5 \\
      \midrule
      Contor & 0 & \at{4}{2} & 0 & \at{4}{3} & \at{4}{1} & \at{4}{1} \\
      \bottomrule
    \end{tabular}
    \end{center}}
  \begin{remarca}<5->{Rezultat}
    Fiecare valoare se scrie de câte ori apare: $[1, 1, 3, 3, 3, 4, 5]$.
  \end{remarca}
  \beginnote{Vizualizare pas cu pas: \href{https://visualgo.net/en/sorting}{visualgo.net/en/sorting}}
\end{frame}

\begin{frame}[fragile]{Counting Sort \textemdash{} implementare}
  \begin{lstlisting}
const int MAXV = 1000000;
int cnt[MAXV + 1];                    // contoarele, initial 0

void countingSort(int v[], int n) {
    int maxVal = 0;
    for (int i = 0; i < n; i++) {
        cnt[v[i]]++;                  // numaram aparitiile
        if (v[i] > maxVal) maxVal = v[i];
    }
    int idx = 0;                      // reconstruim sirul sortat
    for (int val = 0; val <= maxVal; val++) {
        for (int i = 0; i < cnt[val]; i++) {
            v[idx++] = val;
        }
        cnt[val] = 0;                 // contoarele raman 0 pentru apelul urmator
    }
}
  \end{lstlisting}
  \begin{itemize}
    \item Complexitate $O(n + k)$: $n$ pentru numărare, $k$ pentru parcurgerea valorilor posibile
  \end{itemize}
\end{frame}

\begin{frame}{Counting Sort \textemdash{} limitări}
  \begin{intrebare}
    În ce situații Counting Sort nu este potrivit?
  \end{intrebare}
  \begin{block}<2->{Când funcționează}
    \begin{itemize}
      \item Valori întregi nenegative
      \item $k$ mic (de exemplu note $1$\textendash{}$10$, vârste $0$\textendash{}$150$, valori $\leq 10^6$)
    \end{itemize}
  \end{block}
  \begin{block}<3->{Când nu funcționează}
    \begin{itemize}
      \item Valori negative (se adună o constantă tuturor valorilor) sau reale (trebuie transformate mai întâi)
      \item $k$ foarte mare: valori până la $10^9$ ar necesita un tablou de aproximativ $4$ GB
    \end{itemize}
  \end{block}
\end{frame}

% ================================================================
\section{Comparație}

\begin{frame}{Comparația algoritmilor de sortare}
  \begin{center}
  \begin{tabular}{llll}
    \toprule
    Algoritm & Timp & Spațiu suplimentar & Condiții \\
    \midrule
    Bubble Sort    & $O(n^2)$      & $O(1)$      & orice tip comparabil \\
    Selection Sort & $O(n^2)$      & $O(1)$      & orice tip comparabil \\
    Insertion Sort & $O(n^2)$      & $O(1)$      & orice tip comparabil \\
    STL sort       & $O(n \log n)$ & $O(\log n)$ & orice tip comparabil \\
    Counting Sort  & $O(n + k)$    & $O(k)$      & întregi în $[0, k]$ \\
    \bottomrule
  \end{tabular}
  \end{center}
  \begin{intrebare}<2->
    Ce algoritm este potrivit pentru $n = 10^5$ valori până la $10^9$?
  \end{intrebare}
  \uncover<3->{STL sort: metodele pătratice fac circa $10^{10}$ operații, iar Counting Sort ar necesita un tablou de $10^9$ elemente.}
  \begin{itemize}
    \item<4-> În practică se folosește \textbf{STL sort}, dacă nu există restricții speciale
    \item<4-> Counting Sort este mai rapid decât STL sort când $k = O(n)$
    \item<4-> Bubble Sort este util didactic, dar se evită în concursuri
  \end{itemize}
\end{frame}

% ================================================================
\section{Recapitulare}

\begin{frame}[fragile]{Recapitulare}
  \begin{center}
  \begin{tabular}{ll}
    \toprule
    Operație & Cod C++ \\
    \midrule
    Sortare crescătoare    & \lstinline|sort(v, v + n)| \\
    Sortare descrescătoare & \lstinline|sort(v, v + n, greater<int>())| \\
    Sortare a pozițiilor $l..r$ & \lstinline[keepspaces=true]|sort(v + l, v + r + 1)| \\
    Comparator propriu     & \lstinline|sort(v, v + n, cmp)| \\
    \bottomrule
  \end{tabular}
  \end{center}
  \begin{itemize}
    \item \textbf{Bubble, Selection, Insertion Sort}: $O(n^2)$; Bubble și Insertion fac exact atâtea interschimbări câte inversiuni există
    \item \textbf{STL sort}: $O(n \log n)$, soluția implicită; comparatorul nu trebuie să întoarcă \lstinline|true| pentru elemente egale
    \item \textbf{Counting Sort}: $O(n + k)$, doar pentru întregi mici
  \end{itemize}
\end{frame}

\begin{frame}{Probleme de antrenament}
  \begin{block}{Ușoară}
    \begin{enumerate}
      \item \href{https://code.algolymp.com/problems/1081}{\textbf{Even Sort}}
            \textcolor{gray}{(AlgoCode \#1081)}
    \end{enumerate}
  \end{block}
  \begin{block}{Medii}
    \begin{enumerate}
      \item \href{https://code.algolymp.com/problems/1077}{\textbf{K Sort}}
            \textcolor{gray}{(AlgoCode \#1077)}
      \item \href{https://code.algolymp.com/problems/1076}{\textbf{Sort Primes}}
            \textcolor{gray}{(AlgoCode \#1076)}
    \end{enumerate}
  \end{block}
  \begin{block}{Grea}
    \begin{enumerate}
      \item \href{https://code.algolymp.com/problems/1082}{\textbf{Array Task}}
            \textcolor{gray}{(AlgoCode \#1082)}
    \end{enumerate}
  \end{block}
\end{frame}

\begin{frame}{Resurse}
  \begin{block}{Materiale}
    \begin{enumerate}
      \item \href{https://www.pbinfo.ro/articole/5587/sortarea-tablourilor}{\textbf{pbinfo.ro}}
            \textemdash{} Sortarea tablourilor (RO)
      \item \href{https://www.pbinfo.ro/articole/5563/sortarea-prin-selectie}{\textbf{pbinfo.ro}}
            \textemdash{} Sortare prin selecție (RO)
      \item \href{https://www.pbinfo.ro/articole/23702/standard-template-library-stl}{\textbf{pbinfo.ro}}
            \textemdash{} STL în română
      \item \href{https://cp-algorithms.com/sequences/bubble-sort.html}{\textbf{cp-algorithms.com}}
            \textemdash{} Bubble Sort (EN)
      \item \href{https://en.cppreference.com/w/cpp/algorithm/sort}{\textbf{cppreference.com}}
            \textemdash{} std::sort (EN)
      \item \href{https://en.wikipedia.org/wiki/Quicksort}{\textbf{wikipedia.org}}
            \textemdash{} Quicksort și \href{https://en.wikipedia.org/wiki/Merge_sort}{Merge Sort}, algoritmi $O(n \log n)$ (EN)
    \end{enumerate}
  \end{block}
  \begin{block}{Vizualizări}
    \begin{enumerate}
      \item \href{https://visualgo.net/en/sorting}{\textbf{VisuAlgo}}
            \textemdash{} animații pas cu pas pentru Bubble, Selection, Insertion, Merge, Quick, Counting și Radix Sort, cu date introduse de utilizator (EN)
      \item \href{https://www.toptal.com/developers/sorting-algorithms}{\textbf{Toptal}}
            \textemdash{} opt algoritmi comparați pe date aleatoare, aproape sortate, inversate și cu puține valori distincte (EN)
      \item \href{https://www.cs.usfca.edu/~galles/visualization/ComparisonSort.html}{\textbf{USFCA}}
            \textemdash{} Comparison Sorting Visualization (EN)
      \item \href{https://panthema.net/2013/sound-of-sorting/}{\textbf{Sound of Sorting}}
            \textemdash{} vizualizare și sunet pentru operațiile de sortare (EN)
      \item \href{https://www.youtube.com/watch?v=kPRA0W1kECg}{\textbf{YouTube}}
            \textemdash{} 15 Sorting Algorithms in 6 Minutes (EN)
    \end{enumerate}
  \end{block}
\end{frame}

\end{document}
"
% ================================================================
\documentclass[9pt]{beamer}
\usepackage{theme-sorting}

% Mod interactiv: întrebarea apare prima, apoi soluția se dezvăluie pe bucăți
% (\item<2->, \uncover<3->{...}, \begin{remarca}<4->{Concluzie}); hint-ul se
% folosește doar unde elevii chiar trebuie împinși să se gândească.
% Totul se afișează cumulativ: ce apare la pasul 2 rămâne vizibil după.
\newenvironment<>{intrebare}{%
  \begin{uncoverenv}#1%
  \begin{block}{Întrebare}}{\end{block}\end{uncoverenv}}
% Hint și remarcă: text inline cu etichetă emerald, fără cutie
\newenvironment<>{hint}{%
  \begin{uncoverenv}#1%
  \par\textcolor{igMarine}{\textbf{Hint:}}\ }{\par\end{uncoverenv}}
\newenvironment<>{remarca}[1]{%
  \begin{uncoverenv}#2%
  \par\smallskip\textcolor{igMarine}{\textbf{#1:}}\ }{\par\end{uncoverenv}}

% Celulă de tabel vizibilă de la pasul #1
\newcommand{\at}[2]{\uncover<#1->{#2}}

% Diagrame cu tablouri: celulă normală (N), comparată / interschimbată (H),
% la locul definitiv (S), minim sau element inserat (M)
\tikzset{
  cN/.style={draw=igNavy!55,fill=white,minimum size=0.62cm,inner sep=0pt,font=\small},
  cH/.style={cN,fill=igBlue!22,draw=igBlue,thick},
  cS/.style={cN,fill=igSlate!26},
  cM/.style={cN,fill=igOchre!32,draw=igOchre,thick}
}
% \arow{y}{valoare/stil, ...}{text}  (5 elemente)
\newcommand{\arow}[3]{%
  \foreach \v/\t [count=\i from 0] in {#2} {\node[c\t] at (\i*0.68,#1) {\v};}%
  \node[anchor=west,font=\footnotesize] at (5*0.68-0.25,#1) {#3};}

\title{Algoritmi de sortare}
\subtitle{Bubble $\cdot$ Selection $\cdot$ Insertion $\cdot$ STL $\cdot$ Counting Sort}
\author[A. M. Ciobanu]{Andrei Mihai Ciobanu}
\institute[Algolymp]{Algolymp Summer School}
\date{}
\titlegraphic{\includegraphics[width=3cm]{../../../resources/algolymp-logo.png}}

% ================================================================
\begin{document}

\igtitleslide

% ================================================================
\section{Problema sortării}

\begin{frame}{Rolul sortării}
  Sortarea este una dintre cele mai frecvente operații în programare.
  \begin{itemize}
    \item Căutarea rapidă a unui element necesită un șir sortat (tehnică studiată în lecția următoare)
    \item Multe probleme devin simple după sortare: al $k$-lea cel mai mare element, perechi cu suma dată, eliminarea duplicatelor
    \item Date reale: clasamente, liste de note, ierarhii
  \end{itemize}
  \begin{block}{Problema sortării}
    Dat un vector $v[0..n-1]$, să se rearanjeze elementele astfel încât
    $v[0] \leq v[1] \leq \ldots \leq v[n-1]$.
  \end{block}
  \begin{intrebare}<2->
    Cum se verifică dacă un șir de $n$ numere conține două elemente egale, dacă șirul este nesortat, respectiv sortat?
  \end{intrebare}
  \uncover<3->{Nesortat: se compară toate perechile, deci circa $\frac{n(n-1)}{2}$ comparații. Sortat: elementele egale sunt adiacente, deci este suficientă o singură parcurgere, cu $n-1$ comparații.}
\end{frame}

\begin{frame}{Inversiuni}
  O \textbf{inversiune} este o pereche de poziții $i < j$ cu $v[i] > v[j]$.
  \begin{intrebare}
    Câte inversiuni are șirul $[5,\ 3,\ 8,\ 1,\ 4]$?
  \end{intrebare}
  \begin{hint}<2->
    Se numără, pentru fiecare element, valorile mai mici aflate în dreapta lui.
  \end{hint}
  \begin{itemize}
    \item<3-> $5$: $(5,3),\ (5,1),\ (5,4)$ \quad $\rightarrow$ 3 inversiuni
    \item<3-> $3$: $(3,1)$ \quad $\rightarrow$ 1 inversiune
    \item<3-> $8$: $(8,1),\ (8,4)$ \quad $\rightarrow$ 2 inversiuni
    \item<3-> $1$ și $4$: nicio inversiune \quad $\Rightarrow$ \textbf{6 inversiuni} în total
  \end{itemize}
  \begin{remarca}<4->{Legătura cu sortarea}
    Un șir este sortat dacă și numai dacă nu are inversiuni. Interschimbarea a două elemente \textbf{adiacente} aflate în ordine greșită elimină exact o inversiune, iar numărul maxim de inversiuni este $\frac{n(n-1)}{2}$ (șir descrescător).
  \end{remarca}
\end{frame}

% ================================================================
\section{Bubble Sort}

\begin{frame}{Bubble Sort \textemdash{} idee}
  Se compară perechile de elemente adiacente, iar elementele aflate în ordine greșită se interschimbă.
  \begin{intrebare}
    Ce se întâmplă cu cel mai mare element după o parcurgere completă, de la stânga la dreapta, a șirului $[5,\ 3,\ 8,\ 1,\ 4]$?
  \end{intrebare}
  \begin{center}
  \begin{tikzpicture}
    \uncover<2->{\arow{0}{5/N,3/N,8/N,1/N,4/N}{șirul inițial}}
    \uncover<3->{\arow{-0.7}{3/H,5/H,8/N,1/N,4/N}{$5 > 3$: se interschimbă}}
    \uncover<4->{\arow{-1.4}{3/N,5/H,8/H,1/N,4/N}{$5 < 8$: rămân pe loc}}
    \uncover<5->{\arow{-2.1}{3/N,5/N,1/H,8/H,4/N}{$8 > 1$: se interschimbă}}
    \uncover<6->{\arow{-2.8}{3/N,5/N,1/N,4/H,8/H}{$8 > 4$: se interschimbă}}
    \uncover<7->{\arow{-3.5}{3/N,5/N,1/N,4/N,8/S}{sfârșitul pasului 1}}
  \end{tikzpicture}
  \end{center}
  \begin{remarca}<7->{Concluzie}
    După pasul $1$, cel mai mare element se află pe ultima poziție, la locul său definitiv.
  \end{remarca}
  \beginnote{Vizualizare pas cu pas: \href{https://visualgo.net/en/sorting}{visualgo.net/en/sorting}}
\end{frame}

\begin{frame}{Bubble Sort \textemdash{} pașii următori}
  După pasul $i$, ultimele $i$ elemente sunt la locul lor definitiv, deci pasul $i+1$ nu mai ajunge la ele.
  \begin{center}
  \begin{tikzpicture}
    \arow{0}{3/N,5/N,1/N,4/N,8/S}{după pasul 1 \ (3 interschimbări)}
    \uncover<2->{\arow{-0.7}{3/N,1/N,4/N,5/S,8/S}{după pasul 2 \ (2 interschimbări)}}
    \uncover<3->{\arow{-1.4}{1/S,3/S,4/S,5/S,8/S}{după pasul 3 \ (1 interschimbare)}}
  \end{tikzpicture}
  \end{center}
  \begin{intrebare}<4->
    Câte interschimbări a făcut algoritmul în total și ce legătură are numărul cu inversiunile șirului inițial?
  \end{intrebare}
  \uncover<5->{$3 + 2 + 1 = 6$ interschimbări, adică exact cele $6$ inversiuni ale șirului $[5, 3, 8, 1, 4]$.}
  \begin{itemize}
    \item<6-> Pasul $4$ nu mai modifică nimic: pentru $n$ elemente sunt necesari cel mult $n-1$ pași
  \end{itemize}
\end{frame}

\begin{frame}[fragile]{Bubble Sort \textemdash{} implementare}
  \begin{lstlisting}
void bubbleSort(int v[], int n) {
    for (int i = 0; i < n - 1; i++) {
        for (int j = 0; j < n - 1 - i; j++) {
            if (v[j] > v[j + 1]) {
                swap(v[j], v[j + 1]);
            }
        }
    }
}
  \end{lstlisting}
  \begin{itemize}
    \item Bucla exterioară: $n-1$ pași
    \item Bucla interioară: la pasul $i$ se fac $n-1-i$ comparații
    \item Total comparații: $(n-1) + (n-2) + \ldots + 1 = \frac{n(n-1)}{2}$
  \end{itemize}
  \begin{intrebare}<2->
    De ce bucla interioară se oprește la $n-1-i$, nu la $n-1$?
  \end{intrebare}
  \uncover<3->{Ultimele $i$ elemente sunt deja la locul lor, deci nu mai este nevoie să fie comparate.}
  \beginnote{swap \textemdash{} funcție din STL, interschimbă valorile a două variabile}
\end{frame}

\begin{frame}[fragile]{Bubble Sort \textemdash{} optimizare și complexitate}
  \begin{intrebare}
    Ce complexitate are algoritmul de mai sus pe un șir deja sortat?
  \end{intrebare}
  \uncover<2->{
    $O(n^2)$: comparațiile se fac oricum. Algoritmul se poate opri imediat ce o parcurgere completă nu mai găsește nicio inversare.}
  \begin{uncoverenv}<3->
    \begin{lstlisting}
void bubbleSort(int v[], int n) {
    bool ok = true;            // exista inversiuni
    while (ok) {
        ok = false;            // presupunem ca nu mai sunt
        for (int j = 0; j < n - 1; j++) {
            if (v[j] > v[j + 1]) {
                swap(v[j], v[j + 1]);
                ok = true;     // am gasit o inversiune
            }
        }
    }
}
    \end{lstlisting}
  \end{uncoverenv}
  \begin{itemize}
    \item<4-> \textbf{Best case} (șir sortat): $O(n)$ \textemdash{} o singură parcurgere, nicio inversare
    \item<4-> \textbf{Worst case} (șir sortat invers) și \textbf{average case}: $O(n^2)$
  \end{itemize}
\end{frame}

% ================================================================
\section{Selection Sort}

\begin{frame}{Selection Sort \textemdash{} idee}
  La fiecare pas se găsește minimul din partea nesortată și se aduce pe prima poziție a acesteia.
  \begin{intrebare}
    Cum se sortează șirul $[5,\ 3,\ 8,\ 1,\ 4]$ prin alegerea repetată a minimului?
  \end{intrebare}
  \begin{center}
  \begin{tikzpicture}
    \uncover<2->{\arow{0}{5/H,3/N,8/N,1/M,4/N}{minim $1$ $\rightarrow$ se interschimbă cu poziția 0}}
    \uncover<3->{\arow{-0.7}{1/S,3/M,8/N,5/N,4/N}{minim $3$, deja la locul lui}}
    \uncover<4->{\arow{-1.4}{1/S,3/S,8/H,5/N,4/M}{minim $4$ $\rightarrow$ se interschimbă cu poziția 2}}
    \uncover<5->{\arow{-2.1}{1/S,3/S,4/S,5/M,8/N}{minim $5$, deja la locul lui}}
    \uncover<6->{\arow{-2.8}{1/S,3/S,4/S,5/S,8/S}{șir sortat}}
  \end{tikzpicture}
  \end{center}
  \begin{remarca}<6->{Concluzie}
    După pasul $i$, primele $i+1$ elemente sunt la locul lor definitiv.
  \end{remarca}
  \beginnote{Vizualizare pas cu pas: \href{https://visualgo.net/en/sorting}{visualgo.net/en/sorting}}
\end{frame}

\begin{frame}[fragile]{Selection Sort \textemdash{} implementare}
  \begin{lstlisting}
void selectionSort(int v[], int n) {
    for (int i = 0; i < n - 1; i++) {
        int pozMin = i;                  // pozitia minimului din v[i..n-1]
        for (int j = i + 1; j < n; j++) {
            if (v[j] < v[pozMin]) {
                pozMin = j;
            }
        }
        swap(v[i], v[pozMin]);
    }
}
  \end{lstlisting}
  \begin{itemize}
    \item \textbf{Best, worst și average case}: $O(n^2)$ \textemdash{} mereu se fac $\frac{n(n-1)}{2}$ comparații
    \item Se fac cel mult $n-1$ interschimbări (câte una pe pas), mai puține decât la Bubble Sort
  \end{itemize}
  \begin{intrebare}<2->
    Ce complexitate are algoritmul pe un șir deja sortat?
  \end{intrebare}
  \uncover<3->{Tot $O(n^2)$: minimul se caută oricum în toată partea nesortată.}
  \beginnote{swap \textemdash{} funcție din STL, interschimbă valorile a două variabile}
\end{frame}

% ================================================================
\section{Insertion Sort}

\begin{frame}{Insertion Sort \textemdash{} idee}
  Se construiește un prefix sortat, iar fiecare element nou este inserat la locul potrivit în acest prefix.
  \begin{intrebare}
    Cum se extinde un prefix sortat cu încă un element, fără a sorta din nou tot prefixul?
  \end{intrebare}
  \begin{center}
  \begin{tikzpicture}
    \uncover<2->{\arow{0}{5/S,3/M,8/N,1/N,4/N}{se inserează $3$ $\rightarrow$ $[3, 5]$}}
    \uncover<3->{\arow{-0.7}{3/S,5/S,8/M,1/N,4/N}{se inserează $8$ \textemdash{} rămâne pe loc}}
    \uncover<4->{\arow{-1.4}{3/S,5/S,8/S,1/M,4/N}{se inserează $1$ $\rightarrow$ ajunge la început}}
    \uncover<5->{\arow{-2.1}{1/S,3/S,5/S,8/S,4/M}{se inserează $4$ $\rightarrow$ între $3$ și $5$}}
    \uncover<6->{\arow{-2.8}{1/S,3/S,4/S,5/S,8/S}{șir sortat}}
  \end{tikzpicture}
  \end{center}
  \begin{remarca}<6->{Concluzie}
    Elementul inserat este mutat spre stânga cât timp este mai mic decât vecinul său.
  \end{remarca}
  \beginnote{Vizualizare pas cu pas: \href{https://visualgo.net/en/sorting}{visualgo.net/en/sorting}}
\end{frame}

\begin{frame}[fragile]{Insertion Sort \textemdash{} implementare}
  \begin{lstlisting}
void insertionSort(int v[], int n) {
    for (int i = 1; i < n; i++) {
        int j = i;
        while (j > 0 && v[j - 1] > v[j]) {
            swap(v[j - 1], v[j]);
            j--;
        }
    }
}
  \end{lstlisting}
  \begin{itemize}
    \item \textbf{Best case} (șir sortat): $O(n)$ \textemdash{} nicio deplasare
    \item \textbf{Worst case} (șir sortat invers) și \textbf{average case}: $O(n^2)$
  \end{itemize}
  \begin{intrebare}<2->
    De ce condiția buclei conține $j > 0$?
  \end{intrebare}
  \uncover<3->{Fără ea, pentru $j = 0$ s-ar accesa $v[-1]$, adică în afara tabloului.}
  \beginnote{swap \textemdash{} funcție din STL, interschimbă valorile a două variabile}
\end{frame}

\begin{frame}{Cele trei metode pătratice}
  \begin{intrebare}
    Care dintre cele trei metode face cele mai puține interschimbări și care se comportă bine pe un șir aproape sortat?
  \end{intrebare}
  \uncover<2->{
  \begin{center}
  \begin{tabular}{lccc}
    \toprule
    Algoritm & Best case & Worst case & Interschimbări \\
    \midrule
    Bubble Sort (optimizat) & $O(n)$   & $O(n^2)$ & \at{2}{egal cu numărul de inversiuni} \\
    Selection Sort          & $O(n^2)$ & $O(n^2)$ & \at{2}{cel mult $n-1$} \\
    Insertion Sort          & $O(n)$   & $O(n^2)$ & \at{2}{egal cu numărul de inversiuni} \\
    \bottomrule
  \end{tabular}
  \end{center}}
  \begin{itemize}
    \item<3-> \textbf{Selection Sort} face cele mai puține interschimbări
    \item<3-> \textbf{Insertion Sort} este eficient pe șiruri aproape sortate: timpul depinde de numărul de inversiuni
  \end{itemize}
\end{frame}

% ================================================================
\section{STL sort}

\begin{frame}[fragile]{STL sort \textemdash{} utilizare de bază}
  C++ oferă funcția \lstinline|sort| din header-ul \lstinline|<algorithm>|, cu complexitate $O(n \log n)$.
  \begin{intrebare}
    Cum se sortează un tablou $v$ cu $n$ elemente crescător, respectiv descrescător?
  \end{intrebare}
  \begin{uncoverenv}<2->
    \begin{lstlisting}
#include <algorithm>
using namespace std;

int v[] = {5, 3, 8, 1, 4};
int n = 5;

sort(v, v + n);                    // crescator: 1 3 4 5 8
sort(v, v + n, greater<int>());    // descrescator: 8 5 4 3 1
sort(v + l, v + r + 1);            // doar pozitiile l..r

vector<int> w = {5, 3, 8, 1, 4};
sort(w.begin(), w.end());          // pentru vector
    \end{lstlisting}
  \end{uncoverenv}
  \beginnote{vector, begin, end, greater \textemdash{} din STL}
\end{frame}

\begin{frame}[fragile]{STL sort \textemdash{} comparatori}
  Criteriul de sortare poate fi schimbat printr-o \textbf{funcție de comparare}, care întoarce \lstinline|true| dacă primul argument trebuie să apară înaintea celui de-al doilea.
  \begin{intrebare}
    Cum se sortează numerele după valoarea absolută?
  \end{intrebare}
  \begin{uncoverenv}<2->
    \begin{lstlisting}
bool cmpAbs(int a, int b) {
    return abs(a) < abs(b);
}

sort(v, v + n, cmpAbs);
    \end{lstlisting}
  \end{uncoverenv}
  \begin{uncoverenv}<3->
    Pentru perechi: crescător după primul element, iar la egalitate descrescător după al doilea.
    \begin{lstlisting}
bool cmpPair(pair<int, int> a, pair<int, int> b) {
    if (a.first != b.first)
        return a.first < b.first;
    return a.second > b.second;
}

sort(pv, pv + m, cmpPair);   // pv: tablou de pair<int, int>
    \end{lstlisting}
  \end{uncoverenv}
  \beginnote{pair, abs \textemdash{} din STL}
\end{frame}

\begin{frame}[fragile]{STL sort \textemdash{} reguli pentru comparatori}
  \begin{intrebare}
    Ce poate face \lstinline|sort| dacă un comparator întoarce \lstinline|true| și pentru două elemente egale, de exemplu \lstinline|return a <= b;|?
  \end{intrebare}
  \uncover<2->{Comportament nedefinit: programul poate da rezultate greșite sau se poate opri cu eroare (acces în afara tabloului).}
  \begin{block}<3->{Reguli pentru un comparator corect}
    \begin{itemize}
      \item \lstinline|cmp(a, a)| trebuie să întoarcă \lstinline|false|: se folosesc \lstinline|<| și \lstinline|>|, niciodată \lstinline|<=| sau \lstinline|>=|
      \item dacă \lstinline|cmp(a, b)| este \lstinline|true|, atunci \lstinline|cmp(b, a)| trebuie să fie \lstinline|false|
      \item relația trebuie să fie tranzitivă: din \lstinline|cmp(a, b)| și \lstinline|cmp(b, c)| rezultă \lstinline|cmp(a, c)|
    \end{itemize}
  \end{block}
  \begin{itemize}
    \item<4-> \lstinline|sort| nu garantează ordinea relativă a elementelor egale; \lstinline|stable_sort| o păstrează (sortare \textbf{stabilă})
  \end{itemize}
\end{frame}

% ================================================================
\section{Counting Sort}

\begin{frame}{Counting Sort \textemdash{} idee}
  Elementele nu se compară între ele: se numără de câte ori apare fiecare valoare, apoi șirul se reconstruiește.
  \begin{intrebare}
    Cum se poate sorta șirul $v = [3,\ 1,\ 4,\ 1,\ 5,\ 3,\ 3]$ cu valori în intervalul $[0, 5]$, fără nicio comparație între elemente?
  \end{intrebare}
  \begin{hint}<2->
    Valorile sunt numere mici: pot fi folosite ca indici într-un tablou de contoare.
  \end{hint}
  \uncover<3->{
    \begin{center}
    \begin{tabular}{ccccccc}
      \toprule
      Valoare & 0 & 1 & 2 & 3 & 4 & 5 \\
      \midrule
      Contor & 0 & \at{4}{2} & 0 & \at{4}{3} & \at{4}{1} & \at{4}{1} \\
      \bottomrule
    \end{tabular}
    \end{center}}
  \begin{remarca}<5->{Rezultat}
    Fiecare valoare se scrie de câte ori apare: $[1, 1, 3, 3, 3, 4, 5]$.
  \end{remarca}
  \beginnote{Vizualizare pas cu pas: \href{https://visualgo.net/en/sorting}{visualgo.net/en/sorting}}
\end{frame}

\begin{frame}[fragile]{Counting Sort \textemdash{} implementare}
  \begin{lstlisting}
const int MAXV = 1000000;
int cnt[MAXV + 1];                    // contoarele, initial 0

void countingSort(int v[], int n) {
    int maxVal = 0;
    for (int i = 0; i < n; i++) {
        cnt[v[i]]++;                  // numaram aparitiile
        if (v[i] > maxVal) maxVal = v[i];
    }
    int idx = 0;                      // reconstruim sirul sortat
    for (int val = 0; val <= maxVal; val++) {
        for (int i = 0; i < cnt[val]; i++) {
            v[idx++] = val;
        }
        cnt[val] = 0;                 // contoarele raman 0 pentru apelul urmator
    }
}
  \end{lstlisting}
  \begin{itemize}
    \item Complexitate $O(n + k)$: $n$ pentru numărare, $k$ pentru parcurgerea valorilor posibile
  \end{itemize}
\end{frame}

\begin{frame}{Counting Sort \textemdash{} limitări}
  \begin{intrebare}
    În ce situații Counting Sort nu este potrivit?
  \end{intrebare}
  \begin{block}<2->{Când funcționează}
    \begin{itemize}
      \item Valori întregi nenegative
      \item $k$ mic (de exemplu note $1$\textendash{}$10$, vârste $0$\textendash{}$150$, valori $\leq 10^6$)
    \end{itemize}
  \end{block}
  \begin{block}<3->{Când nu funcționează}
    \begin{itemize}
      \item Valori negative (se adună o constantă tuturor valorilor) sau reale (trebuie transformate mai întâi)
      \item $k$ foarte mare: valori până la $10^9$ ar necesita un tablou de aproximativ $4$ GB
    \end{itemize}
  \end{block}
\end{frame}

% ================================================================
\section{Comparație}

\begin{frame}{Comparația algoritmilor de sortare}
  \begin{center}
  \begin{tabular}{llll}
    \toprule
    Algoritm & Timp & Spațiu suplimentar & Condiții \\
    \midrule
    Bubble Sort    & $O(n^2)$      & $O(1)$      & orice tip comparabil \\
    Selection Sort & $O(n^2)$      & $O(1)$      & orice tip comparabil \\
    Insertion Sort & $O(n^2)$      & $O(1)$      & orice tip comparabil \\
    STL sort       & $O(n \log n)$ & $O(\log n)$ & orice tip comparabil \\
    Counting Sort  & $O(n + k)$    & $O(k)$      & întregi în $[0, k]$ \\
    \bottomrule
  \end{tabular}
  \end{center}
  \begin{intrebare}<2->
    Ce algoritm este potrivit pentru $n = 10^5$ valori până la $10^9$?
  \end{intrebare}
  \uncover<3->{STL sort: metodele pătratice fac circa $10^{10}$ operații, iar Counting Sort ar necesita un tablou de $10^9$ elemente.}
  \begin{itemize}
    \item<4-> În practică se folosește \textbf{STL sort}, dacă nu există restricții speciale
    \item<4-> Counting Sort este mai rapid decât STL sort când $k = O(n)$
    \item<4-> Bubble Sort este util didactic, dar se evită în concursuri
  \end{itemize}
\end{frame}

% ================================================================
\section{Recapitulare}

\begin{frame}[fragile]{Recapitulare}
  \begin{center}
  \begin{tabular}{ll}
    \toprule
    Operație & Cod C++ \\
    \midrule
    Sortare crescătoare    & \lstinline|sort(v, v + n)| \\
    Sortare descrescătoare & \lstinline|sort(v, v + n, greater<int>())| \\
    Sortare a pozițiilor $l..r$ & \lstinline[keepspaces=true]|sort(v + l, v + r + 1)| \\
    Comparator propriu     & \lstinline|sort(v, v + n, cmp)| \\
    \bottomrule
  \end{tabular}
  \end{center}
  \begin{itemize}
    \item \textbf{Bubble, Selection, Insertion Sort}: $O(n^2)$; Bubble și Insertion fac exact atâtea interschimbări câte inversiuni există
    \item \textbf{STL sort}: $O(n \log n)$, soluția implicită; comparatorul nu trebuie să întoarcă \lstinline|true| pentru elemente egale
    \item \textbf{Counting Sort}: $O(n + k)$, doar pentru întregi mici
  \end{itemize}
\end{frame}

\begin{frame}{Probleme de antrenament}
  \begin{block}{Ușoară}
    \begin{enumerate}
      \item \href{https://code.algolymp.com/problems/1081}{\textbf{Even Sort}}
            \textcolor{gray}{(AlgoCode \#1081)}
    \end{enumerate}
  \end{block}
  \begin{block}{Medii}
    \begin{enumerate}
      \item \href{https://code.algolymp.com/problems/1077}{\textbf{K Sort}}
            \textcolor{gray}{(AlgoCode \#1077)}
      \item \href{https://code.algolymp.com/problems/1076}{\textbf{Sort Primes}}
            \textcolor{gray}{(AlgoCode \#1076)}
    \end{enumerate}
  \end{block}
  \begin{block}{Grea}
    \begin{enumerate}
      \item \href{https://code.algolymp.com/problems/1082}{\textbf{Array Task}}
            \textcolor{gray}{(AlgoCode \#1082)}
    \end{enumerate}
  \end{block}
\end{frame}

\begin{frame}{Resurse}
  \begin{block}{Materiale}
    \begin{enumerate}
      \item \href{https://www.pbinfo.ro/articole/5587/sortarea-tablourilor}{\textbf{pbinfo.ro}}
            \textemdash{} Sortarea tablourilor (RO)
      \item \href{https://www.pbinfo.ro/articole/5563/sortarea-prin-selectie}{\textbf{pbinfo.ro}}
            \textemdash{} Sortare prin selecție (RO)
      \item \href{https://www.pbinfo.ro/articole/23702/standard-template-library-stl}{\textbf{pbinfo.ro}}
            \textemdash{} STL în română
      \item \href{https://cp-algorithms.com/sequences/bubble-sort.html}{\textbf{cp-algorithms.com}}
            \textemdash{} Bubble Sort (EN)
      \item \href{https://en.cppreference.com/w/cpp/algorithm/sort}{\textbf{cppreference.com}}
            \textemdash{} std::sort (EN)
      \item \href{https://en.wikipedia.org/wiki/Quicksort}{\textbf{wikipedia.org}}
            \textemdash{} Quicksort și \href{https://en.wikipedia.org/wiki/Merge_sort}{Merge Sort}, algoritmi $O(n \log n)$ (EN)
    \end{enumerate}
  \end{block}
  \begin{block}{Vizualizări}
    \begin{enumerate}
      \item \href{https://visualgo.net/en/sorting}{\textbf{VisuAlgo}}
            \textemdash{} animații pas cu pas pentru Bubble, Selection, Insertion, Merge, Quick, Counting și Radix Sort, cu date introduse de utilizator (EN)
      \item \href{https://www.toptal.com/developers/sorting-algorithms}{\textbf{Toptal}}
            \textemdash{} opt algoritmi comparați pe date aleatoare, aproape sortate, inversate și cu puține valori distincte (EN)
      \item \href{https://www.cs.usfca.edu/~galles/visualization/ComparisonSort.html}{\textbf{USFCA}}
            \textemdash{} Comparison Sorting Visualization (EN)
      \item \href{https://panthema.net/2013/sound-of-sorting/}{\textbf{Sound of Sorting}}
            \textemdash{} vizualizare și sunet pentru operațiile de sortare (EN)
      \item \href{https://www.youtube.com/watch?v=kPRA0W1kECg}{\textbf{YouTube}}
            \textemdash{} 15 Sorting Algorithms in 6 Minutes (EN)
    \end{enumerate}
  \end{block}
\end{frame}

\end{document}
"
% ================================================================
\documentclass[9pt]{beamer}
\usepackage{theme-sorting}

% Mod interactiv: întrebarea apare prima, apoi soluția se dezvăluie pe bucăți
% (\item<2->, \uncover<3->{...}, \begin{remarca}<4->{Concluzie}); hint-ul se
% folosește doar unde elevii chiar trebuie împinși să se gândească.
% Totul se afișează cumulativ: ce apare la pasul 2 rămâne vizibil după.
\newenvironment<>{intrebare}{%
  \begin{uncoverenv}#1%
  \begin{block}{Întrebare}}{\end{block}\end{uncoverenv}}
% Hint și remarcă: text inline cu etichetă emerald, fără cutie
\newenvironment<>{hint}{%
  \begin{uncoverenv}#1%
  \par\textcolor{igMarine}{\textbf{Hint:}}\ }{\par\end{uncoverenv}}
\newenvironment<>{remarca}[1]{%
  \begin{uncoverenv}#2%
  \par\smallskip\textcolor{igMarine}{\textbf{#1:}}\ }{\par\end{uncoverenv}}

% Celulă de tabel vizibilă de la pasul #1
\newcommand{\at}[2]{\uncover<#1->{#2}}

% Diagrame cu tablouri: celulă normală (N), comparată / interschimbată (H),
% la locul definitiv (S), minim sau element inserat (M)
\tikzset{
  cN/.style={draw=igNavy!55,fill=white,minimum size=0.62cm,inner sep=0pt,font=\small},
  cH/.style={cN,fill=igBlue!22,draw=igBlue,thick},
  cS/.style={cN,fill=igSlate!26},
  cM/.style={cN,fill=igOchre!32,draw=igOchre,thick}
}
% \arow{y}{valoare/stil, ...}{text}  (5 elemente)
\newcommand{\arow}[3]{%
  \foreach \v/\t [count=\i from 0] in {#2} {\node[c\t] at (\i*0.68,#1) {\v};}%
  \node[anchor=west,font=\footnotesize] at (5*0.68-0.25,#1) {#3};}

\title{Algoritmi de sortare}
\subtitle{Bubble $\cdot$ Selection $\cdot$ Insertion $\cdot$ STL $\cdot$ Counting Sort}
\author[A. M. Ciobanu]{Andrei Mihai Ciobanu}
\institute[Algolymp]{Algolymp Summer School}
\date{}
\titlegraphic{\includegraphics[width=3cm]{../../../resources/algolymp-logo.png}}

% ================================================================
\begin{document}

\igtitleslide

% ================================================================
\section{Problema sortării}

\begin{frame}{Rolul sortării}
  Sortarea este una dintre cele mai frecvente operații în programare.
  \begin{itemize}
    \item Căutarea rapidă a unui element necesită un șir sortat (tehnică studiată în lecția următoare)
    \item Multe probleme devin simple după sortare: al $k$-lea cel mai mare element, perechi cu suma dată, eliminarea duplicatelor
    \item Date reale: clasamente, liste de note, ierarhii
  \end{itemize}
  \begin{block}{Problema sortării}
    Dat un vector $v[0..n-1]$, să se rearanjeze elementele astfel încât
    $v[0] \leq v[1] \leq \ldots \leq v[n-1]$.
  \end{block}
  \begin{intrebare}<2->
    Cum se verifică dacă un șir de $n$ numere conține două elemente egale, dacă șirul este nesortat, respectiv sortat?
  \end{intrebare}
  \uncover<3->{Nesortat: se compară toate perechile, deci circa $\frac{n(n-1)}{2}$ comparații. Sortat: elementele egale sunt adiacente, deci este suficientă o singură parcurgere, cu $n-1$ comparații.}
\end{frame}

\begin{frame}{Inversiuni}
  O \textbf{inversiune} este o pereche de poziții $i < j$ cu $v[i] > v[j]$.
  \begin{intrebare}
    Câte inversiuni are șirul $[5,\ 3,\ 8,\ 1,\ 4]$?
  \end{intrebare}
  \begin{hint}<2->
    Se numără, pentru fiecare element, valorile mai mici aflate în dreapta lui.
  \end{hint}
  \begin{itemize}
    \item<3-> $5$: $(5,3),\ (5,1),\ (5,4)$ \quad $\rightarrow$ 3 inversiuni
    \item<3-> $3$: $(3,1)$ \quad $\rightarrow$ 1 inversiune
    \item<3-> $8$: $(8,1),\ (8,4)$ \quad $\rightarrow$ 2 inversiuni
    \item<3-> $1$ și $4$: nicio inversiune \quad $\Rightarrow$ \textbf{6 inversiuni} în total
  \end{itemize}
  \begin{remarca}<4->{Legătura cu sortarea}
    Un șir este sortat dacă și numai dacă nu are inversiuni. Interschimbarea a două elemente \textbf{adiacente} aflate în ordine greșită elimină exact o inversiune, iar numărul maxim de inversiuni este $\frac{n(n-1)}{2}$ (șir descrescător).
  \end{remarca}
\end{frame}

% ================================================================
\section{Bubble Sort}

\begin{frame}{Bubble Sort \textemdash{} idee}
  Se compară perechile de elemente adiacente, iar elementele aflate în ordine greșită se interschimbă.
  \begin{intrebare}
    Ce se întâmplă cu cel mai mare element după o parcurgere completă, de la stânga la dreapta, a șirului $[5,\ 3,\ 8,\ 1,\ 4]$?
  \end{intrebare}
  \begin{center}
  \begin{tikzpicture}
    \uncover<2->{\arow{0}{5/N,3/N,8/N,1/N,4/N}{șirul inițial}}
    \uncover<3->{\arow{-0.7}{3/H,5/H,8/N,1/N,4/N}{$5 > 3$: se interschimbă}}
    \uncover<4->{\arow{-1.4}{3/N,5/H,8/H,1/N,4/N}{$5 < 8$: rămân pe loc}}
    \uncover<5->{\arow{-2.1}{3/N,5/N,1/H,8/H,4/N}{$8 > 1$: se interschimbă}}
    \uncover<6->{\arow{-2.8}{3/N,5/N,1/N,4/H,8/H}{$8 > 4$: se interschimbă}}
    \uncover<7->{\arow{-3.5}{3/N,5/N,1/N,4/N,8/S}{sfârșitul pasului 1}}
  \end{tikzpicture}
  \end{center}
  \begin{remarca}<7->{Concluzie}
    După pasul $1$, cel mai mare element se află pe ultima poziție, la locul său definitiv.
  \end{remarca}
  \beginnote{Vizualizare pas cu pas: \href{https://visualgo.net/en/sorting}{visualgo.net/en/sorting}}
\end{frame}

\begin{frame}{Bubble Sort \textemdash{} pașii următori}
  După pasul $i$, ultimele $i$ elemente sunt la locul lor definitiv, deci pasul $i+1$ nu mai ajunge la ele.
  \begin{center}
  \begin{tikzpicture}
    \arow{0}{3/N,5/N,1/N,4/N,8/S}{după pasul 1 \ (3 interschimbări)}
    \uncover<2->{\arow{-0.7}{3/N,1/N,4/N,5/S,8/S}{după pasul 2 \ (2 interschimbări)}}
    \uncover<3->{\arow{-1.4}{1/S,3/S,4/S,5/S,8/S}{după pasul 3 \ (1 interschimbare)}}
  \end{tikzpicture}
  \end{center}
  \begin{intrebare}<4->
    Câte interschimbări a făcut algoritmul în total și ce legătură are numărul cu inversiunile șirului inițial?
  \end{intrebare}
  \uncover<5->{$3 + 2 + 1 = 6$ interschimbări, adică exact cele $6$ inversiuni ale șirului $[5, 3, 8, 1, 4]$.}
  \begin{itemize}
    \item<6-> Pasul $4$ nu mai modifică nimic: pentru $n$ elemente sunt necesari cel mult $n-1$ pași
  \end{itemize}
\end{frame}

\begin{frame}[fragile]{Bubble Sort \textemdash{} implementare}
  \begin{lstlisting}
void bubbleSort(int v[], int n) {
    for (int i = 0; i < n - 1; i++) {
        for (int j = 0; j < n - 1 - i; j++) {
            if (v[j] > v[j + 1]) {
                swap(v[j], v[j + 1]);
            }
        }
    }
}
  \end{lstlisting}
  \begin{itemize}
    \item Bucla exterioară: $n-1$ pași
    \item Bucla interioară: la pasul $i$ se fac $n-1-i$ comparații
    \item Total comparații: $(n-1) + (n-2) + \ldots + 1 = \frac{n(n-1)}{2}$
  \end{itemize}
  \begin{intrebare}<2->
    De ce bucla interioară se oprește la $n-1-i$, nu la $n-1$?
  \end{intrebare}
  \uncover<3->{Ultimele $i$ elemente sunt deja la locul lor, deci nu mai este nevoie să fie comparate.}
  \beginnote{swap \textemdash{} funcție din STL, interschimbă valorile a două variabile}
\end{frame}

\begin{frame}[fragile]{Bubble Sort \textemdash{} optimizare și complexitate}
  \begin{intrebare}
    Ce complexitate are algoritmul de mai sus pe un șir deja sortat?
  \end{intrebare}
  \uncover<2->{
    $O(n^2)$: comparațiile se fac oricum. Algoritmul se poate opri imediat ce o parcurgere completă nu mai găsește nicio inversare.}
  \begin{uncoverenv}<3->
    \begin{lstlisting}
void bubbleSort(int v[], int n) {
    bool ok = true;            // exista inversiuni
    while (ok) {
        ok = false;            // presupunem ca nu mai sunt
        for (int j = 0; j < n - 1; j++) {
            if (v[j] > v[j + 1]) {
                swap(v[j], v[j + 1]);
                ok = true;     // am gasit o inversiune
            }
        }
    }
}
    \end{lstlisting}
  \end{uncoverenv}
  \begin{itemize}
    \item<4-> \textbf{Best case} (șir sortat): $O(n)$ \textemdash{} o singură parcurgere, nicio inversare
    \item<4-> \textbf{Worst case} (șir sortat invers) și \textbf{average case}: $O(n^2)$
  \end{itemize}
\end{frame}

% ================================================================
\section{Selection Sort}

\begin{frame}{Selection Sort \textemdash{} idee}
  La fiecare pas se găsește minimul din partea nesortată și se aduce pe prima poziție a acesteia.
  \begin{intrebare}
    Cum se sortează șirul $[5,\ 3,\ 8,\ 1,\ 4]$ prin alegerea repetată a minimului?
  \end{intrebare}
  \begin{center}
  \begin{tikzpicture}
    \uncover<2->{\arow{0}{5/H,3/N,8/N,1/M,4/N}{minim $1$ $\rightarrow$ se interschimbă cu poziția 0}}
    \uncover<3->{\arow{-0.7}{1/S,3/M,8/N,5/N,4/N}{minim $3$, deja la locul lui}}
    \uncover<4->{\arow{-1.4}{1/S,3/S,8/H,5/N,4/M}{minim $4$ $\rightarrow$ se interschimbă cu poziția 2}}
    \uncover<5->{\arow{-2.1}{1/S,3/S,4/S,5/M,8/N}{minim $5$, deja la locul lui}}
    \uncover<6->{\arow{-2.8}{1/S,3/S,4/S,5/S,8/S}{șir sortat}}
  \end{tikzpicture}
  \end{center}
  \begin{remarca}<6->{Concluzie}
    După pasul $i$, primele $i+1$ elemente sunt la locul lor definitiv.
  \end{remarca}
  \beginnote{Vizualizare pas cu pas: \href{https://visualgo.net/en/sorting}{visualgo.net/en/sorting}}
\end{frame}

\begin{frame}[fragile]{Selection Sort \textemdash{} implementare}
  \begin{lstlisting}
void selectionSort(int v[], int n) {
    for (int i = 0; i < n - 1; i++) {
        int pozMin = i;                  // pozitia minimului din v[i..n-1]
        for (int j = i + 1; j < n; j++) {
            if (v[j] < v[pozMin]) {
                pozMin = j;
            }
        }
        swap(v[i], v[pozMin]);
    }
}
  \end{lstlisting}
  \begin{itemize}
    \item \textbf{Best, worst și average case}: $O(n^2)$ \textemdash{} mereu se fac $\frac{n(n-1)}{2}$ comparații
    \item Se fac cel mult $n-1$ interschimbări (câte una pe pas), mai puține decât la Bubble Sort
  \end{itemize}
  \begin{intrebare}<2->
    Ce complexitate are algoritmul pe un șir deja sortat?
  \end{intrebare}
  \uncover<3->{Tot $O(n^2)$: minimul se caută oricum în toată partea nesortată.}
  \beginnote{swap \textemdash{} funcție din STL, interschimbă valorile a două variabile}
\end{frame}

% ================================================================
\section{Insertion Sort}

\begin{frame}{Insertion Sort \textemdash{} idee}
  Se construiește un prefix sortat, iar fiecare element nou este inserat la locul potrivit în acest prefix.
  \begin{intrebare}
    Cum se extinde un prefix sortat cu încă un element, fără a sorta din nou tot prefixul?
  \end{intrebare}
  \begin{center}
  \begin{tikzpicture}
    \uncover<2->{\arow{0}{5/S,3/M,8/N,1/N,4/N}{se inserează $3$ $\rightarrow$ $[3, 5]$}}
    \uncover<3->{\arow{-0.7}{3/S,5/S,8/M,1/N,4/N}{se inserează $8$ \textemdash{} rămâne pe loc}}
    \uncover<4->{\arow{-1.4}{3/S,5/S,8/S,1/M,4/N}{se inserează $1$ $\rightarrow$ ajunge la început}}
    \uncover<5->{\arow{-2.1}{1/S,3/S,5/S,8/S,4/M}{se inserează $4$ $\rightarrow$ între $3$ și $5$}}
    \uncover<6->{\arow{-2.8}{1/S,3/S,4/S,5/S,8/S}{șir sortat}}
  \end{tikzpicture}
  \end{center}
  \begin{remarca}<6->{Concluzie}
    Elementul inserat este mutat spre stânga cât timp este mai mic decât vecinul său.
  \end{remarca}
  \beginnote{Vizualizare pas cu pas: \href{https://visualgo.net/en/sorting}{visualgo.net/en/sorting}}
\end{frame}

\begin{frame}[fragile]{Insertion Sort \textemdash{} implementare}
  \begin{lstlisting}
void insertionSort(int v[], int n) {
    for (int i = 1; i < n; i++) {
        int j = i;
        while (j > 0 && v[j - 1] > v[j]) {
            swap(v[j - 1], v[j]);
            j--;
        }
    }
}
  \end{lstlisting}
  \begin{itemize}
    \item \textbf{Best case} (șir sortat): $O(n)$ \textemdash{} nicio deplasare
    \item \textbf{Worst case} (șir sortat invers) și \textbf{average case}: $O(n^2)$
  \end{itemize}
  \begin{intrebare}<2->
    De ce condiția buclei conține $j > 0$?
  \end{intrebare}
  \uncover<3->{Fără ea, pentru $j = 0$ s-ar accesa $v[-1]$, adică în afara tabloului.}
  \beginnote{swap \textemdash{} funcție din STL, interschimbă valorile a două variabile}
\end{frame}

\begin{frame}{Cele trei metode pătratice}
  \begin{intrebare}
    Care dintre cele trei metode face cele mai puține interschimbări și care se comportă bine pe un șir aproape sortat?
  \end{intrebare}
  \uncover<2->{
  \begin{center}
  \begin{tabular}{lccc}
    \toprule
    Algoritm & Best case & Worst case & Interschimbări \\
    \midrule
    Bubble Sort (optimizat) & $O(n)$   & $O(n^2)$ & \at{2}{egal cu numărul de inversiuni} \\
    Selection Sort          & $O(n^2)$ & $O(n^2)$ & \at{2}{cel mult $n-1$} \\
    Insertion Sort          & $O(n)$   & $O(n^2)$ & \at{2}{egal cu numărul de inversiuni} \\
    \bottomrule
  \end{tabular}
  \end{center}}
  \begin{itemize}
    \item<3-> \textbf{Selection Sort} face cele mai puține interschimbări
    \item<3-> \textbf{Insertion Sort} este eficient pe șiruri aproape sortate: timpul depinde de numărul de inversiuni
  \end{itemize}
\end{frame}

% ================================================================
\section{STL sort}

\begin{frame}[fragile]{STL sort \textemdash{} utilizare de bază}
  C++ oferă funcția \lstinline|sort| din header-ul \lstinline|<algorithm>|, cu complexitate $O(n \log n)$.
  \begin{intrebare}
    Cum se sortează un tablou $v$ cu $n$ elemente crescător, respectiv descrescător?
  \end{intrebare}
  \begin{uncoverenv}<2->
    \begin{lstlisting}
#include <algorithm>
using namespace std;

int v[] = {5, 3, 8, 1, 4};
int n = 5;

sort(v, v + n);                    // crescator: 1 3 4 5 8
sort(v, v + n, greater<int>());    // descrescator: 8 5 4 3 1
sort(v + l, v + r + 1);            // doar pozitiile l..r

vector<int> w = {5, 3, 8, 1, 4};
sort(w.begin(), w.end());          // pentru vector
    \end{lstlisting}
  \end{uncoverenv}
  \beginnote{vector, begin, end, greater \textemdash{} din STL}
\end{frame}

\begin{frame}[fragile]{STL sort \textemdash{} comparatori}
  Criteriul de sortare poate fi schimbat printr-o \textbf{funcție de comparare}, care întoarce \lstinline|true| dacă primul argument trebuie să apară înaintea celui de-al doilea.
  \begin{intrebare}
    Cum se sortează numerele după valoarea absolută?
  \end{intrebare}
  \begin{uncoverenv}<2->
    \begin{lstlisting}
bool cmpAbs(int a, int b) {
    return abs(a) < abs(b);
}

sort(v, v + n, cmpAbs);
    \end{lstlisting}
  \end{uncoverenv}
  \begin{uncoverenv}<3->
    Pentru perechi: crescător după primul element, iar la egalitate descrescător după al doilea.
    \begin{lstlisting}
bool cmpPair(pair<int, int> a, pair<int, int> b) {
    if (a.first != b.first)
        return a.first < b.first;
    return a.second > b.second;
}

sort(pv, pv + m, cmpPair);   // pv: tablou de pair<int, int>
    \end{lstlisting}
  \end{uncoverenv}
  \beginnote{pair, abs \textemdash{} din STL}
\end{frame}

\begin{frame}[fragile]{STL sort \textemdash{} reguli pentru comparatori}
  \begin{intrebare}
    Ce poate face \lstinline|sort| dacă un comparator întoarce \lstinline|true| și pentru două elemente egale, de exemplu \lstinline|return a <= b;|?
  \end{intrebare}
  \uncover<2->{Comportament nedefinit: programul poate da rezultate greșite sau se poate opri cu eroare (acces în afara tabloului).}
  \begin{block}<3->{Reguli pentru un comparator corect}
    \begin{itemize}
      \item \lstinline|cmp(a, a)| trebuie să întoarcă \lstinline|false|: se folosesc \lstinline|<| și \lstinline|>|, niciodată \lstinline|<=| sau \lstinline|>=|
      \item dacă \lstinline|cmp(a, b)| este \lstinline|true|, atunci \lstinline|cmp(b, a)| trebuie să fie \lstinline|false|
      \item relația trebuie să fie tranzitivă: din \lstinline|cmp(a, b)| și \lstinline|cmp(b, c)| rezultă \lstinline|cmp(a, c)|
    \end{itemize}
  \end{block}
  \begin{itemize}
    \item<4-> \lstinline|sort| nu garantează ordinea relativă a elementelor egale; \lstinline|stable_sort| o păstrează (sortare \textbf{stabilă})
  \end{itemize}
\end{frame}

% ================================================================
\section{Counting Sort}

\begin{frame}{Counting Sort \textemdash{} idee}
  Elementele nu se compară între ele: se numără de câte ori apare fiecare valoare, apoi șirul se reconstruiește.
  \begin{intrebare}
    Cum se poate sorta șirul $v = [3,\ 1,\ 4,\ 1,\ 5,\ 3,\ 3]$ cu valori în intervalul $[0, 5]$, fără nicio comparație între elemente?
  \end{intrebare}
  \begin{hint}<2->
    Valorile sunt numere mici: pot fi folosite ca indici într-un tablou de contoare.
  \end{hint}
  \uncover<3->{
    \begin{center}
    \begin{tabular}{ccccccc}
      \toprule
      Valoare & 0 & 1 & 2 & 3 & 4 & 5 \\
      \midrule
      Contor & 0 & \at{4}{2} & 0 & \at{4}{3} & \at{4}{1} & \at{4}{1} \\
      \bottomrule
    \end{tabular}
    \end{center}}
  \begin{remarca}<5->{Rezultat}
    Fiecare valoare se scrie de câte ori apare: $[1, 1, 3, 3, 3, 4, 5]$.
  \end{remarca}
  \beginnote{Vizualizare pas cu pas: \href{https://visualgo.net/en/sorting}{visualgo.net/en/sorting}}
\end{frame}

\begin{frame}[fragile]{Counting Sort \textemdash{} implementare}
  \begin{lstlisting}
const int MAXV = 1000000;
int cnt[MAXV + 1];                    // contoarele, initial 0

void countingSort(int v[], int n) {
    int maxVal = 0;
    for (int i = 0; i < n; i++) {
        cnt[v[i]]++;                  // numaram aparitiile
        if (v[i] > maxVal) maxVal = v[i];
    }
    int idx = 0;                      // reconstruim sirul sortat
    for (int val = 0; val <= maxVal; val++) {
        for (int i = 0; i < cnt[val]; i++) {
            v[idx++] = val;
        }
        cnt[val] = 0;                 // contoarele raman 0 pentru apelul urmator
    }
}
  \end{lstlisting}
  \begin{itemize}
    \item Complexitate $O(n + k)$: $n$ pentru numărare, $k$ pentru parcurgerea valorilor posibile
  \end{itemize}
\end{frame}

\begin{frame}{Counting Sort \textemdash{} limitări}
  \begin{intrebare}
    În ce situații Counting Sort nu este potrivit?
  \end{intrebare}
  \begin{block}<2->{Când funcționează}
    \begin{itemize}
      \item Valori întregi nenegative
      \item $k$ mic (de exemplu note $1$\textendash{}$10$, vârste $0$\textendash{}$150$, valori $\leq 10^6$)
    \end{itemize}
  \end{block}
  \begin{block}<3->{Când nu funcționează}
    \begin{itemize}
      \item Valori negative (se adună o constantă tuturor valorilor) sau reale (trebuie transformate mai întâi)
      \item $k$ foarte mare: valori până la $10^9$ ar necesita un tablou de aproximativ $4$ GB
    \end{itemize}
  \end{block}
\end{frame}

% ================================================================
\section{Comparație}

\begin{frame}{Comparația algoritmilor de sortare}
  \begin{center}
  \begin{tabular}{llll}
    \toprule
    Algoritm & Timp & Spațiu suplimentar & Condiții \\
    \midrule
    Bubble Sort    & $O(n^2)$      & $O(1)$      & orice tip comparabil \\
    Selection Sort & $O(n^2)$      & $O(1)$      & orice tip comparabil \\
    Insertion Sort & $O(n^2)$      & $O(1)$      & orice tip comparabil \\
    STL sort       & $O(n \log n)$ & $O(\log n)$ & orice tip comparabil \\
    Counting Sort  & $O(n + k)$    & $O(k)$      & întregi în $[0, k]$ \\
    \bottomrule
  \end{tabular}
  \end{center}
  \begin{intrebare}<2->
    Ce algoritm este potrivit pentru $n = 10^5$ valori până la $10^9$?
  \end{intrebare}
  \uncover<3->{STL sort: metodele pătratice fac circa $10^{10}$ operații, iar Counting Sort ar necesita un tablou de $10^9$ elemente.}
  \begin{itemize}
    \item<4-> În practică se folosește \textbf{STL sort}, dacă nu există restricții speciale
    \item<4-> Counting Sort este mai rapid decât STL sort când $k = O(n)$
    \item<4-> Bubble Sort este util didactic, dar se evită în concursuri
  \end{itemize}
\end{frame}

% ================================================================
\section{Recapitulare}

\begin{frame}[fragile]{Recapitulare}
  \begin{center}
  \begin{tabular}{ll}
    \toprule
    Operație & Cod C++ \\
    \midrule
    Sortare crescătoare    & \lstinline|sort(v, v + n)| \\
    Sortare descrescătoare & \lstinline|sort(v, v + n, greater<int>())| \\
    Sortare a pozițiilor $l..r$ & \lstinline[keepspaces=true]|sort(v + l, v + r + 1)| \\
    Comparator propriu     & \lstinline|sort(v, v + n, cmp)| \\
    \bottomrule
  \end{tabular}
  \end{center}
  \begin{itemize}
    \item \textbf{Bubble, Selection, Insertion Sort}: $O(n^2)$; Bubble și Insertion fac exact atâtea interschimbări câte inversiuni există
    \item \textbf{STL sort}: $O(n \log n)$, soluția implicită; comparatorul nu trebuie să întoarcă \lstinline|true| pentru elemente egale
    \item \textbf{Counting Sort}: $O(n + k)$, doar pentru întregi mici
  \end{itemize}
\end{frame}

\begin{frame}{Probleme de antrenament}
  \begin{block}{Ușoară}
    \begin{enumerate}
      \item \href{https://code.algolymp.com/problems/1081}{\textbf{Even Sort}}
            \textcolor{gray}{(AlgoCode \#1081)}
    \end{enumerate}
  \end{block}
  \begin{block}{Medii}
    \begin{enumerate}
      \item \href{https://code.algolymp.com/problems/1077}{\textbf{K Sort}}
            \textcolor{gray}{(AlgoCode \#1077)}
      \item \href{https://code.algolymp.com/problems/1076}{\textbf{Sort Primes}}
            \textcolor{gray}{(AlgoCode \#1076)}
    \end{enumerate}
  \end{block}
  \begin{block}{Grea}
    \begin{enumerate}
      \item \href{https://code.algolymp.com/problems/1082}{\textbf{Array Task}}
            \textcolor{gray}{(AlgoCode \#1082)}
    \end{enumerate}
  \end{block}
\end{frame}

\begin{frame}{Resurse}
  \begin{block}{Materiale}
    \begin{enumerate}
      \item \href{https://www.pbinfo.ro/articole/5587/sortarea-tablourilor}{\textbf{pbinfo.ro}}
            \textemdash{} Sortarea tablourilor (RO)
      \item \href{https://www.pbinfo.ro/articole/5563/sortarea-prin-selectie}{\textbf{pbinfo.ro}}
            \textemdash{} Sortare prin selecție (RO)
      \item \href{https://www.pbinfo.ro/articole/23702/standard-template-library-stl}{\textbf{pbinfo.ro}}
            \textemdash{} STL în română
      \item \href{https://cp-algorithms.com/sequences/bubble-sort.html}{\textbf{cp-algorithms.com}}
            \textemdash{} Bubble Sort (EN)
      \item \href{https://en.cppreference.com/w/cpp/algorithm/sort}{\textbf{cppreference.com}}
            \textemdash{} std::sort (EN)
      \item \href{https://en.wikipedia.org/wiki/Quicksort}{\textbf{wikipedia.org}}
            \textemdash{} Quicksort și \href{https://en.wikipedia.org/wiki/Merge_sort}{Merge Sort}, algoritmi $O(n \log n)$ (EN)
    \end{enumerate}
  \end{block}
  \begin{block}{Vizualizări}
    \begin{enumerate}
      \item \href{https://visualgo.net/en/sorting}{\textbf{VisuAlgo}}
            \textemdash{} animații pas cu pas pentru Bubble, Selection, Insertion, Merge, Quick, Counting și Radix Sort, cu date introduse de utilizator (EN)
      \item \href{https://www.toptal.com/developers/sorting-algorithms}{\textbf{Toptal}}
            \textemdash{} opt algoritmi comparați pe date aleatoare, aproape sortate, inversate și cu puține valori distincte (EN)
      \item \href{https://www.cs.usfca.edu/~galles/visualization/ComparisonSort.html}{\textbf{USFCA}}
            \textemdash{} Comparison Sorting Visualization (EN)
      \item \href{https://panthema.net/2013/sound-of-sorting/}{\textbf{Sound of Sorting}}
            \textemdash{} vizualizare și sunet pentru operațiile de sortare (EN)
      \item \href{https://www.youtube.com/watch?v=kPRA0W1kECg}{\textbf{YouTube}}
            \textemdash{} 15 Sorting Algorithms in 6 Minutes (EN)
    \end{enumerate}
  \end{block}
\end{frame}

\end{document}
